%% ============================================================
%%  SUPPLEMENTARY MATERIAL B
%%  HypatiaX: Noise Robustness, Sample Complexity, and Convergence
%%  of Hybrid LLM–Neural Symbolic Regression on a Custom 30-Equation Benchmark
%%
%%  Companion to: jmlr_paper_main.tex
%%  Sources: six-method noiseless comparison (clean report)  +
%%           noise/sample-complexity sweep (v2, 16 March 2026)
%%  Rewritten & restructured: March 2026
%% ============================================================

\documentclass[twoside,11pt]{article}

%% ── JMLR house style ─────────────────────────────────────────────────────────
\usepackage{jmlr2e}

%% ── Core packages ────────────────────────────────────────────────────────────
\usepackage{amsmath,amssymb}
\usepackage{booktabs}
\usepackage{multirow}
\usepackage{makecell}
\usepackage{xcolor}
\usepackage{colortbl}
\usepackage{array}
\usepackage{longtable}% generated per-case tables (llm_detailed, defi_detailed)
\usepackage{caption}
\usepackage{subcaption}
\usepackage[expansion=true]{microtype}
\usepackage{rotating}
\usepackage{algorithm}
\usepackage{algorithmic}
\usepackage{graphicx}
\usepackage{float}
\usepackage{pgfplots}
\pgfplotsset{compat=1.18}
\usepackage{cleveref}
\usepackage{verbatim}
\usepackage{listings}

%% ── Code listings ────────────────────────────────────────────────────────────
\definecolor{codeblue}{RGB}{0,90,170}
\definecolor{codegray}{RGB}{100,100,100}
\definecolor{codegreen}{RGB}{0,128,64}
\definecolor{codered}{RGB}{180,30,30}
\definecolor{backgray}{RGB}{248,248,248}

\lstset{
  language=Python,
  basicstyle=\ttfamily\small,
  keywordstyle=\color{codeblue}\bfseries,
  commentstyle=\color{codegray}\itshape,
  stringstyle=\color{codegreen},
  backgroundcolor=\color{backgray},
  frame=single,
  framerule=0.4pt,
  rulecolor=\color{black!20},
  breaklines=true,
  breakatwhitespace=true,
  numbers=left,
  numberstyle=\tiny\color{codegray},
  numbersep=8pt,
  xleftmargin=12pt,
  tabsize=4,
  showstringspaces=false,
  captionpos=b
}

%% ── Colours ──────────────────────────────────────────────────────────────────
\definecolor{passgreen}{RGB}{22,163,74}
\definecolor{passgreenBg}{RGB}{220,252,231}
\definecolor{failred}{RGB}{220,38,38}
\definecolor{failredBg}{RGB}{254,226,226}
\definecolor{warnAmber}{RGB}{217,119,6}
\definecolor{warnAmberBg}{RGB}{254,243,199}
\definecolor{headnavyBg}{RGB}{219,234,254}
\definecolor{m3blue}{RGB}{31,119,180}
\definecolor{m4red}{RGB}{214,39,40}
\definecolor{tabhead}{RGB}{220,230,245}

%% ── Column types ─────────────────────────────────────────────────────────────
\newcolumntype{C}[1]{>{\centering\arraybackslash}p{#1}}
\newcolumntype{L}[1]{>{\raggedright\arraybackslash}p{#1}}
\newcolumntype{R}[1]{>{\raggedleft\arraybackslash}p{#1}}

%% ── Method macros ────────────────────────────────────────────────────────────
\newcommand{\PureLLM}{\textsc{PureLLM}}
\newcommand{\INN}{\textsc{ImpNN}}
\newcommand{\EHD}{\textsc{EHSDeFi}}
\newcommand{\HSL}{\textsc{HLLMNN}}
\newcommand{\SEL}{\textsc{SymLLM}}
\newcommand{\HDS}{\textsc{HDSv50\_2}}
\newcommand{\Mthree}{{\color{m3blue}\textbf{\EHD}}}
\newcommand{\Mfour}{{\color{m4red}\textbf{\HSL}}}

%% ── Metric shorthand ─────────────────────────────────────────────────────────
\newcommand{\Rsq}{R^{2}}
\newcommand{\rmse}{\ensuremath{\mathrm{RMSE}}}
\newcommand{\nrmse}{\ensuremath{\widetilde{\mathrm{RMSE}}}}
\newcommand{\warn}{{\color{warnAmber}\textbf{!}}}

%% ── Theorem-like environments ────────────────────────────────────────────────
%% NOTE: amsthm is deliberately NOT loaded and `example` is NOT
%% redeclared here. jmlr2e.sty already defines proof/theorem/lemma/
%% proposition/remark/corollary/definition/example unconditionally
%% via plain-LaTeX \newtheorem/\newenvironment. Loading amsthm here
%% interferes with jmlr2e's optional-argument counter machinery
%% (symptom: "No counter 'Definition' defined" elsewhere in the
%% submission), and redeclaring `example` collides outright
%% ("Command \theorem already defined"-class error). Neither
%% needed — jmlr2e's `example` environment (shared, continuous
%% numbering, not per-section) already supports the
%% \begin{example}[Title] usage below.

%% ============================================================
\begin{document}

\title{%
  Supplementary Material B\\[6pt]
  HypatiaX: Noise Robustness, Sample Complexity, and Convergence\\
  of Hybrid LLM--Neural Symbolic Regression on a Custom 30-Equation Benchmark
}
\author{Ruperto Pedro Bonet Chaple\\
  \texttt{ruperto.bonet@modelphysmat.com}\\
  Independent Researcher, London, United Kingdom
}
\editor{}
\maketitle

\begin{abstract}
We present an extended evaluation of the HypatiaX symbolic regression
framework on a custom 30-equation benchmark, integrating a six-method
noiseless comparison with a comprehensive noise robustness sweep
($\sigma \in \{0, 0.05, 0.1, 0.5, 1\}\%$) and a
sample complexity sweep
($n \in \{50, 100, 200, 500, 750, 1000\}$), focusing on the two
top-performing core methods: \EHD{} (M3) and \HSL{} (M4).

Following correction of a measurement harness bug in \HSL{} (see
Section~\ref{sec:bugfix}), both methods achieve median $\Rsq = 1.000000$
across all tested conditions.
\EHD{} achieves 100\% recovery at all \emph{noisy} conditions ($\sigma > 0$)
but only 76.7\% (23/30) at $\sigma=0\%$ under the strict noiseless
threshold ($\Rsq \geq 0.999999$; \Cref{tab:rr_noise}).

\HSL{} achieves 100\% recovery at all \emph{noisy} conditions ($\sigma > 0$)
and 100\% (30/30 equations, v2-corrected) under the strict noiseless protocol
($\Rsq > 0.9999$). \HDS{} achieves 70.0\% (21/30 equations) under the same
noiseless protocol.

All results in this document use the v2 corrected harness; v1 figures are
retained in the repository commit history for audit purposes only.
Both reach median $\Rsq=1.000000$ at $n = 50$ samples (noiseless sample-complexity
sweep; noisy $n>50$ runs are not available), though \EHD's
recovery rate dips to 96.7\% at larger $n$ (\Cref{tab:sc_metrics}).
\EHD{} offers exact symbolic interpretability; \HSL{} is $\sim$1.1--1.2$\times$
faster at every noise level .

On this 30-equation custom benchmark (\textbf{not} a Feynman subset ---
see \Cref{tab:overall}), \EHD{} achieves 76.7\% recovery and \HDS{} achieves
70.0\%, alongside AI~Feynman~2.0's published 79.3\% on the unrelated
100-equation Feynman/SRBench suite
.
The AI Feynman figure is shown for general SOTA context only: published
baselines were evaluated on
all 100 Feynman equations under the standardised SRBench protocol; these
figures should not be read as a general state-of-the-art claim.
We additionally correct four methodological issues:
undisclosed hardcoded \PureLLM{} results, scale-biased raw \rmse,
missing Wilcoxon significance tests, and mismatched pass-rate thresholds.
\end{abstract}


\begin{keywords}
symbolic regression, scientific discovery, large language models,
neural networks, equation-recovery benchmark, noise robustness, sample complexity,
hybrid systems, reproducibility
\end{keywords}

%% ============================================================
\section{Introduction}
\label{suppb:intro}

This document is Supplementary Material~B to the main JMLR submission
\emph{HypatiaX: A Hybrid Symbolic-Neural Framework for Extrapolation-Reliable
Analytical Discovery} (\texttt{jmlr\_paper\_main.tex}).

\medskip
\noindent\textbf{Cross-reference note (v5, updated 2026-07-28).}
The main paper has been updated since this supplementary was written, and
again since the v4 cross-reference below was added. Readers are directed
to the following sections for context:
\begin{itemize}
  \item \textbf{The main paper's \emph{Protocol-30 Extrapolation Benchmark ($n=30$)} subsection (Results section)} --- the main paper
    reports the 2026-08-12 results: \textbf{21/30 (70.0\,\%)} for
    HybridDiscoverySystem v50\_2 under the random 80/20 split and
    \textbf{18/30 (60.0\,\%)} under the PCA-directed 40\,\%/60\,\% split, at
    the strict $\Rsq \geq 0.999999$ threshold on the method-reported $\Rsq$ (13/30 on the
    held-out 60\,\% for the PCA split). The same
    runs give PureLLM 29/30 and 29/30, ImprovedNN 0/30 and 3/30,
    EnhancedHybridSystemDeFi 23/30 and 23/30, HybridSystemLLMNN 30/30 and
    30/30 (all routed to the LLM), and SymbolicEngineWithLLM 27/30 and 23/30.
    These are per-equation pass counts on the same 30 equations, read from
    \texttt{benchmark\_results.json} and \texttt{benchmark\_results\_pca\_4060.json}
    rather than from the file behind \Cref{tab:overall}. The counts for
    EnhancedHybridSystemDeFi (23/30, 76.7\%) and HybridDiscoverySystem
    (21/30, 70.0\%) coincide with that table. SymbolicEngineWithLLM is 27/30
    here and 28/30 there; see the tablenote on \Cref{tab:overall}.
  \item \textbf{The main paper's \emph{Nguyen-12 Benchmark (Standard SR Suite)} subsection} --- reports
    PySR-only and HypatiaX on the standard Nguyen-12 suite under an axis-aligned
    extrapolation protocol. \textbf{Figures revised 2026-09-28} and rebuilt from
    the consolidated CI result files: seed~42 gives 10/12 for both systems, and
    across five seeds HypatiaX recovers 51/60 (85.0\,\%) and PySR-only 49/60
    (81.7\,\%) at $\Rsq\ge0.9999$. HypatiaX fails Nguyen-12 on all five seeds
    (LLM gate rejects all candidates; the run aborts before any search).
    The earlier 8/12~P / 7/12~H figures, the 11/12 / 10/12 figure, the
    4/12 caption figure and the Neural MLP (0/12) figure, all cited here in
    earlier versions, are retracted: none reproduces from any current source
    file, and no Neural MLP was ever run for this benchmark. See the main paper
    (the Nguyen-12 table, tablenote~\ddag, and the ``Nguyen-12 Benchmark (Standard SR Suite)'' section). The
    pipeline uses per-equation axis-aligned variable ranges, not PCA-directed
    ordering.
  \item \textbf{Section~7.4 Proposition~1 (Routing Cascade Convergence)} --- formal
    guarantee that Phase~2 warm-starting is monotonically non-degrading,
    complementing the routing improvements described in Supplementary~A.
\end{itemize}
The noise robustness, sample complexity, and convergence analyses in this
document remain valid and unchanged.

Symbolic regression (SR) aims to discover closed-form analytical expressions
from data without prior knowledge of the functional form.  The Feynman SR
Benchmark~\citep{udrescu2020ai} provides a canonical evaluation suite of
100 physics equations (the standard set used under SRBench).

This supplementary integrates two complementary evaluations:
\begin{enumerate}
  \item A \textbf{six-method noiseless comparison} (corrected from the
    companion report) establishing the state-of-the-art context. \textbf{[Note:
    this evaluation uses a custom 30-equation benchmark, not a subset of the
    Feynman suite described above. Its 11 domains are biology, chemistry,
    electrochemistry, electromagnetism, electrostatics, magnetism, mechanics,
    optics, probability, quantum, and thermodynamics; several of these
    (biology, chemistry, electrochemistry, probability) are outside the
    Feynman benchmark's scope. An earlier version of this note listed ten
    domains including mathematics and economics; that list does not match
    the source files. See
    \Cref{sec:noiseless} and the source-file note on \Cref{tab:overall}.]}
  \item A \textbf{sweep evaluation} (v2, post bug-fix, 16 March 2026)
    covering \EHD{} (M3) and \HSL{} (M4) across noise levels and sample
    sizes, with convergence, win-rate, and architecture analysis.
\end{enumerate}

\paragraph{Version note.}
A v1 sweep (15 March 2026) incorrectly reported \HSL{} as having a persistent
catastrophic failure on Newton's gravitational force ($\Rsq \approx 0.42$--$0.87$).
This was caused by a measurement bug in \texttt{evaluate\_llm\_formula}
(Section~\ref{sec:bugfix}), not an architectural limitation.  All results in
this paper are from the v2 rerun unless explicitly labelled v1.

%% ============================================================
\section{Related Work}
\label{suppb:related}

SR has evolved from genetic programming~\citep{schmidt2009distilling} through
physics-inspired decomposition~\citep{udrescu2020ai} to transformer-based
generation~\citep{biggio2021neural,shojaee2023transformer} and deep RL over
expression trees~\citep{petersen2021deep}.  The best-published result on the
noiseless Feynman benchmark prior to this work is AI~Feynman~2.0 at
79.3\%~\citep{udrescu2020aifeynman2}.

Benchmark integrity has received increasing attention.
\citet{memorization2024audit} show that LLMs can reproduce iconic equations
from training data rather than inferring them from observations, inflating
reported pass rates.  The hardcoding audit applied to \PureLLM{} in
Section~\ref{sec:integrity} follows this line of work.

A methodological concern specific to tiny-scale physics equations is that
benchmark harnesses using absolute sum-of-squares thresholds silently
misclassify correct formulas.  The fix described in Section~\ref{sec:bugfix}
is a general contribution to SR benchmarking practice.

%% ============================================================
\section{Methods}
\label{sec:methods}

\subsection{Six-Method Benchmark Suite}
\label{sec:methods:sixmethod}

\begin{table}[ht]
\caption{Methods under evaluation: six-method noiseless benchmark.}
\label{tab:methods}
\centering
\small
\renewcommand{\arraystretch}{1.3}
\begin{tabular}{L{2.1cm} L{1.3cm} L{8.0cm}}
\toprule
\textbf{Method} & \textbf{Type} & \textbf{Description} \\
\midrule
\rowcolor{warnAmberBg}
\PureLLM & Core &
  Zero-shot LLM formula inference.  No numerical fitting.
  \textbf{\warn~11/30 results flagged hardcoded (Section~\ref{sec:integrity}).} \\
\EHD\,(M3) & Core &
  Ensemble-first system (140/150 = 93.3\% of routing decisions; \Cref{tab:arch}) combining LLM-guided generation,
  NN ensemble scoring, and DeepFit optimisation. \\
\HSL\,(M4) & Core &
  LLM-first + NN residual correction (123/150 = 82.0\% LLM-first routing decisions; \Cref{tab:arch});
  $\sim$1.1--1.2$\times$ faster than M3 (\Cref{tab:time_noise}). \\
\INN & Core &
  Pure neural baseline ($1\!\to\!128\!\to\!64\!\to\!32\!\to\!1$, 3 seeds). \\
\SEL & Tools &
  LLM-guided symbolic search with CAS verification. \\
\HDS & Tools &
  HybridDiscoverySystem~v50\_2; 21/30 pass at 75.3\,s mean runtime in the noiseless protocol, against 28/30 at 39.1\,s for \SEL{} . \\
\bottomrule
\end{tabular}
\end{table}

\subsection{Sweep Evaluation Design (M3 vs.\ M4)}
\label{sec:methods:sweep}

\begin{table}[ht]
\centering
\caption{Sweep design.}
\label{tab:sweeps}
\begin{tabular}{L{2.5cm} L{4.3cm} L{2.5cm} C{2.0cm}}
\toprule
\textbf{Sweep} & \textbf{Axis varied} & \textbf{Fixed} & \textbf{Conditions}\\
\midrule
Noise       & $\sigma \in \{0, 0.05, 0.1, 0.5, 1\}\%$\textsuperscript{*} & $n=200$       & 5 levels \\
Sample size & $n \in \{50,100,200,500,750,1000\}$  & $\sigma=5\%$ & 6 sizes  \\
\bottomrule
\multicolumn{4}{l}{\footnotesize\textsuperscript{*}Noise levels as recorded in the sweep source files; see \Cref{tab:r2_noise}.}\\
\end{tabular}
\end{table}

\noindent\textbf{Hyperparameters:} NN seeds = 5, random seed 42,
method timeout = 900\,s, PySR timeout = 1100\,s,
threshold (noise sweep) = 0.999999 at $\sigma=0$, $0.995$ at
$\sigma\in\{0.05,0.1,0.5\}\%$ and $0.99$ at $\sigma=1\%$
(sample-complexity sweep, noiseless: $0.999$ at $n=50$, $0.9995$ at
$n=100$, $0.9999$ at $n\geq200$).
All 30 equations completed with zero timeouts across all 11 conditions.

\subsection{Integrity Annotation: Hardcoded Results}
\label{sec:integrity}

Following~\citet{memorization2024audit}, each \PureLLM{} noiseless row is audited for whether the language model
was called. The \texttt{is\_hardcoded} flag is set when the baseline returns
\texttt{method="pure\_llm\_hardcoded"}, i.e.\ when \texttt{\_try\_hardcoded\_formula} in
\texttt{baseline\_pure\_llm\_defi\_discovery.py} returns a formula without any LLM call. It is a record of
that code path, not a statistical memorisation test. Two kinds of bypass exist: eight equations return the
benchmark's own ground-truth expression (GT in Table~\ref{tab:hardcoded}), and three use a fixed functional form whose parameters are fitted to the
data by least squares (Fit). A twelfth branch (Arrhenius) did not fire. We call the remaining 19 rows \emph{live}.
The same 11 equations are identified in the eleven per-domain \texttt{protocol\_core\_noiseless\_20260812\_*.json} files (by flag), in the consolidated \texttt{benchmark\_results.json} and in \texttt{benchmark\_results\_pca\_4060.json} (by timing, 0.5--1.6\,ms, since those two flat files carry no flag), using \texttt{recount\_hardcoded.py}. The \texttt{170344} file was not part of the recount.
11 of 30 \PureLLM{} noiseless rows bypassed the model (Table~\ref{tab:hardcoded}).

\begin{table}[ht]
\centering
\caption{Equations for which \PureLLM{} bypassed the language model (flagged \texttt{is\_hardcoded=True}).
  GT: the benchmark's ground-truth expression is returned; Fit: fixed form with least-squares parameters fitted to the data.
  $\Rsq$ as recorded in the recount; every bypassed row passes.}
\label{tab:hardcoded}
\small
\begin{tabular}{L{5.4cm} L{3.0cm} C{1.0cm} C{1.7cm}}
\toprule
\textbf{Equation} & \textbf{Domain} & \textbf{Type} & \textbf{$\Rsq$ (NL)} \\
\midrule
Michaelis--Menten enzyme kinetics       & Biology          & Fit & 1.0000000 \\
Logistic growth rate                    & Biology          & Fit & 1.0000000 \\
Allometric scaling law                  & Biology          & Fit & 1.0000000 \\
Nernst equation (electrode potential)   & Electrochemistry & GT  & 1.0000000 \\
Clausius--Mossotti effective field      & Electrostatics   & GT  & 1.0000000 \\
Lorentz force: $F = qvB$                & Magnetism        & GT  & 1.0000000 \\
Gaussian/normal distribution PDF        & Probability      & GT  & 1.0000000 \\
Zeeman energy (electron spin)           & Quantum          & GT  & 1.0000000 \\
Bose--Einstein occupation number        & Quantum          & GT  & 1.0000000 \\
Planck blackbody spectral radiance      & Thermodynamics   & GT  & 1.0000000 \\
Fourier's law of heat conduction        & Thermodynamics   & GT  & 1.0000000 \\
\bottomrule
\end{tabular}
\end{table}

\subsection{Evaluation Metrics}
\label{sec:methods:metrics}

\paragraph{Coefficient of determination ($\Rsq$).}
$\Rsq = 1 - \mathrm{SS}_{\mathrm{res}} / \mathrm{SS}_{\mathrm{tot}}$.

\paragraph{Normalised RMSE ($\nrmse$).}
$\nrmse = \rmse/\sigma_y = \sqrt{1 - \Rsq}$.
Raw \rmse{} is dominated by large-scale equations; median \nrmse{} or $\Rsq$
should be used for cross-equation comparison.

\paragraph{Recovery rate.}
Fraction of equations achieving $\Rsq \geq $ threshold ($\geq 0.9999$
noiseless; $\geq 0.995$ noisy).  Catastrophic = $\Rsq < 0.90$.

\paragraph{Statistical tests.}
Wilcoxon signed-rank tests for all pairwise comparisons, Bonferroni corrected
($\alpha^* = 0.05/15 \approx 0.0033$ for $\binom{6}{2}=15$ pairs).

%% ============================================================
\section{Algorithm}
\label{sec:algorithm}

\begin{algorithm}[ht]
\caption{HypatiaX hybrid symbolic discovery (M3/M4 core loop)}
\label{alg:hds}
\begin{algorithmic}[1]
\REQUIRE Dataset $\mathcal{D}$, LLM proposer $\mathcal{M}$,
         NN ensemble $f_\theta$, threshold $\tau$
\ENSURE  Discovered expression $\hat{e}$
\STATE $d \leftarrow \textsc{Describe}(\mathcal{D})$
\STATE $\{e_1,\dots,e_K\} \leftarrow \mathcal{M}(d, \mathcal{D})$
\FOR{each candidate $e_k$}
  \STATE Fit $\theta^* \leftarrow \arg\min_\theta \sum_i (f_\theta(e_k,x_i)-y_i)^2$
  \STATE Compute $\Rsq_k$ with relative threshold (Section~\ref{sec:bugfix})
  \IF{$e_k$ invalid or numerically unstable}
    \STATE \textbf{discard}
  \ENDIF
\ENDFOR
\STATE \textbf{M3 path:} $\hat{e} \leftarrow \arg\max_{e_k} \Rsq_k$
\STATE \textbf{M4 path:} if $\Rsq_{\text{LLM}} \geq \tau$, return LLM formula;
       else apply NN residual correction
\RETURN $\hat{e}$
\end{algorithmic}
\end{algorithm}

%% ============================================================
\section{Six-Method Noiseless Results}
\label{sec:noiseless}

\begin{table}[ht]
\caption{Six-method aggregate performance, noiseless protocol
  ($\Rsq \geq 0.999999$, 30 equations).
  \warn~\PureLLM{} pass rate includes 11/30 hardcoded results.
  Source: \texttt{protocol\_core\_noiseless\_20260812\_170344.json} for all rows except \PureLLM{}. The \PureLLM{} row is a recount over the eleven per-domain \texttt{protocol\_core\_noiseless\_20260812\_*.json} files: 29/30 over all rows, of which 11 are bypassed rows (\Cref{tab:hardcoded}) that all pass, and 18/19 on the \emph{live} rows where the model was called; the 170344 file reports 30/30 (the one-test difference is Curie's law, which returned no formula in \texttt{benchmark\_results.json} and passes in the 170344 file). This is a custom benchmark, not the Feynman-30 suite (\Cref{suppb:intro}); its 11 domains are biology, chemistry, electrochemistry, electromagnetism, electrostatics, magnetism, mechanics, optics, probability, quantum, and thermodynamics.}
\label{tab:overall}
\centering
\small
\renewcommand{\arraystretch}{1.3}
\begin{tabular}{L{2.4cm} C{1.0cm}
                C{1.8cm} C{1.3cm} C{1.4cm}}
\toprule
\textbf{Method} & \textbf{Type} &
  \makecell{\textbf{Pass Rate}\\$(\Rsq \geq 0.999999)$} &
  \makecell{\textbf{Mean} $\Rsq$\\(valid tests)} &
  \makecell{\textbf{Avg}\\Runtime} \\
\midrule
\rowcolor{warnAmberBg}
\PureLLM\,\warn$^\dagger$ & Core  & 29/30 (96.7\%)\warn & --- &   4.7\,s \\
\rowcolor{passgreenBg}
\EHD\,(M3)      & Core  & 23/30 (76.7\%)      & 0.999998 &   7.6\,s \\
\HDS            & Tools & 21/30 (70.0\%)      & 0.999473 &  75.3\,s \\
\HSL\,(M4)      & Core  & 30/30 (100.0\%)     & 1.000000  &   6.9\,s \\
\SEL\,$^\ddagger$ & Tools & 28/30 (93.3\%)      & 0.999711 &  39.1\,s \\
\rowcolor{failredBg}
\INN            & Core  & 0/30 (0.0\%)        & 0.999222 &   2.5\,s \\
\bottomrule
\multicolumn{5}{l}{\footnotesize $^\dagger$\,Recount: 29/30 over all rows, 18/19 live rows, 11/11 bypassed rows. Mean $\Rsq$
  and runtime were not recomputed; the 4.7\,s average includes bypassed rows, which take under 2\,ms.}\\
\multicolumn{5}{p{8.4cm}}{\footnotesize $^\ddagger$\,Read from the 170344 file. The main paper's random-split table, built from \texttt{benchmark\_results.json}, gives 27/30 for \SEL{} (mean $\Rsq$ 0.999868; 23/30 on the PCA split). The two files also differ for \INN{} (mean 0.999222 here, 0.999323 there; 0/30 in both) and for \PureLLM{} (30/30 in the 170344 file, 29/30 in the recount). Means for \EHD, \HDS{} and \HSL{} agree to six decimals. The \PureLLM{} difference is Curie's law, which returned no formula ($\Rsq=0$) in \texttt{benchmark\_results.json}; the \SEL{} difference is Arrhenius ($\Rsq=0.999921$ there, $0.9999999987$ here, threshold $0.999999$), and \SEL{} is not deterministic across runs; \INN{} differs through seed variation. Neither file supersedes the other.}\\
\end{tabular}
\end{table}

\subsection{Context Comparison with Published SR Systems}
\label{sec:sota}

\begin{table}[ht]
\caption{Performance of this work's methods on the 30-equation custom
  noiseless benchmark (\textbf{not} a Feynman subset --- see
  \Cref{suppb:intro} and the source-file note on \Cref{tab:overall}),
  shown alongside published SR systems' figures on the unrelated,
  standard 100-equation Feynman/SRBench suite for general context only.
  \PureLLM{} excluded (Section~\ref{sec:integrity}).
  \textbf{Note:} the published baselines (AI Feynman, NeSymReS, TPSR, DSR)
  and this work's methods were \textbf{not} evaluated on the same
  equation set; the ``vs AI Feynman'' deltas below are illustrative only
  and should not be read as a head-to-head benchmark result.}
\label{tab:sota}
\centering
\small
\begin{tabular}{L{4.2cm} L{2.0cm} C{1.5cm} L{3.5cm}}
\toprule
\textbf{System} & \textbf{Type} & \textbf{Pass Rate} & \textbf{Note} \\
\midrule
\rowcolor{passgreenBg}
\EHD{} (this work) & Hybrid   & \textbf{76.7\%} & $-2.6$\,pp vs AI Feynman (different sets)
   \\
\rowcolor{passgreenBg}
\HDS{} (this work) & Hybrid   & \textbf{70.0\%} & $-9.3$\,pp vs AI Feynman (different sets)
   \\
\midrule
AI~Feynman~2.0     & Symbolic    & 79.3\% & \citet{udrescu2020aifeynman2} (100-eq, SRBench) \\
NeSymReS           & Neural      & 59.4\% & \citet{biggio2021neural} (100-eq, SRBench) \\
TPSR               & Transformer & 56.0\% & \citet{shojaee2023transformer} (100-eq, SRBench) \\
DSR                & Deep SR     & 32.0\% & \citet{petersen2021deep} (100-eq, SRBench) \\
\bottomrule
\end{tabular}
\end{table}

\subsection{Scale-Invariant RMSE}
\label{sec:noiseless:nrmse}

After replacing raw \rmse{} with $\nrmse = \sqrt{1-\Rsq}$, four of six methods
cluster in mean $\nrmse \in [0.000, 0.021]$; \PureLLM{} is the outlier on mean
($\nrmse_{\text{mean}}=0.1042$) and also on max ($\nrmse_{\max}=2.0063$,
driven by the negative-$\Rsq$ ``Thin Lens'' failure disclosed in
\Cref{tab:overall}).  \PureLLM{} and \HDS{} each have raw \rmse{} values
of \texttt{inf}/\texttt{nan}/sentinel-overflow in the source data for
specific equations (see table note), so their mean raw \rmse{} over all
30 equations is not meaningful; \nrmse{} remains well-defined for every
equation with a finite $\Rsq$ since it is derived directly from $\Rsq$.

\begin{table}[ht]
\centering
\caption{Raw RMSE vs.\ normalised RMSE ($\nrmse=\sqrt{1-\Rsq}$), 30 equations.
Regenerated from the eleven per-domain \texttt{protocol\_core\_noiseless\_20260812\_*.json} files (not the 170344 file that supplies five rows of \Cref{tab:overall}; see its tablenote).}
\label{tab:nrmse}
\small
\begin{tabular}{L{3.0cm} C{2.6cm} C{2.0cm} C{2.0cm}}
\toprule
\textbf{Method} & \textbf{Mean RMSE} & \textbf{Mean NRMSE} & \textbf{Max NRMSE} \\
\midrule
\PureLLM (core) & 0.0876$^\dagger$ & 0.1042 & 2.0063 \\
\EHD     (core) & 0.5672        & 0.0037 & 0.0735 \\
\HDS     (tools)& 2{,}740.0$^\ddagger$ & 0.0199 & 0.1675 \\
\SEL     (tools)& 106.7          & 0.0099 & 0.2091 \\
\INN     (core) & 361.4          & 0.0209 & 0.2419 \\
\HSL     (core) & 0.0002         & 0.0002 & 0.0038 \\
\bottomrule
\multicolumn{4}{p{11cm}}{\footnotesize $^\dagger$\,Two equations
  (``Kinetic Energy,'' ``Logistic Growth'') have raw \rmse{}$=\infty$ in
  the source file, from the catastrophic/truncated failures disclosed in
  \Cref{tab:overall}, and are excluded from this mean. Over the 28
  equations with finite \rmse{}, mean \rmse{}$\,\approx 0.0876$.}\\
\multicolumn{4}{p{11cm}}{\footnotesize $^\ddagger$\,Three equations
  (Coulomb's Law, Lorentz Force, Diffraction) have raw \rmse{}$=$\texttt{nan}
  in the source file despite good $\Rsq$ on the same equations, and one
  (Bernoulli's Equation) reports a $10^{150}$ sentinel/overflow value; all
  four are excluded from this mean. Over the 26 remaining equations, mean
  \rmse{}$\,\approx 2{,}740$. This \rmse{}-specific irregularity is
  unexplained and is noted for the record rather than resolved here.}\\
\end{tabular}
\end{table}

\subsection{Wilcoxon Significance Tests}
\label{sec:noiseless:wilcoxon}

\begin{table}[ht]
\centering
\caption{Wilcoxon signed-rank tests on noiseless $\Rsq$, all pairwise comparisons (two-sided,
         30 equations; $n_{\mathrm{eff}}$ is the number of non-zero paired differences,
         Bonferroni $\alpha^*=0.0033$).
         **$\alpha^*$; *$\alpha=0.05$ only; n.s.~not significant.
         Recomputed from the eleven per-domain noiseless files (the source of the \PureLLM{} recount in \Cref{tab:overall}; its other rows come from the 170344 file); invalid/\texttt{nan} $\Rsq$ values are treated as 0 for the paired test. $n_{\mathrm{eff}}$ is below 30 where methods tie.}
\label{tab:wilcoxon}
\small
\begin{tabular}{llrrrl}
\toprule
\textbf{Method A} & \textbf{Method B} & \textbf{Statistic} & $\boldsymbol{n_{\mathrm{eff}}}$ &
  \textbf{$p$-value} & \textbf{Sig.} \\
\midrule
\PureLLM & \INN  &  99.0 & 30 &  0.0050  & *  \\
\PureLLM & \EHD  & 112.0 & 30 &  0.0120  & *  \\
\PureLLM & \HSL  &   2.0 &  7 &  0.0422  & *  \\
\PureLLM & \SEL  &  87.0 & 21 &  0.3219  & n.s. \\
\PureLLM & \HDS  &  90.0 & 23 &  0.1443  & n.s. \\
\midrule
\INN     & \EHD  &  52.0 & 30 &  0.0001  & ** \\
\INN     & \HSL  &   0.0 & 30 & $<0.001$ & ** \\
\INN     & \SEL  &  78.0 & 30 &  0.0010  & ** \\
\INN     & \HDS  & 146.0 & 30 &  0.0767  & n.s. \\
\midrule
\EHD     & \HSL  &  74.0 & 30 &  0.0007  & ** \\
\EHD     & \SEL  & 154.0 & 30 &  0.1094  & n.s. \\
\EHD     & \HDS  & 196.0 & 30 &  0.4645  & n.s. \\
\midrule
\HSL     & \SEL  &  58.0 & 22 &  0.0262  & *  \\
\HSL     & \HDS  &  43.0 & 24 &  0.0022  & ** \\
\SEL     & \HDS  &  81.0 & 24 &  0.0487  & *  \\
\bottomrule
\end{tabular}

\smallskip
{\footnotesize \textit{Notes.} (i) Pairs with $n_{\mathrm{eff}}=30$ have exact $p$-values; pairs with
$n_{\mathrm{eff}}<30$ use the normal approximation, which is unreliable at $n_{\mathrm{eff}}=7$.
(ii) Three starred results are sensitive to the approximation: with a continuity correction,
\PureLLM{} vs.\ \HSL{} gives $p=0.052$ and \SEL{} vs.\ \HDS{} gives $p=0.050$ (both would be not significant at
$\alpha=0.05$), while \HSL{} vs.\ \SEL{} stays at $p=0.027$. None of the three survives the Bonferroni correction.
(iii) \PureLLM{} rows include the 11 bypassed equations (\Cref{tab:hardcoded}); comparisons involving \PureLLM{}
therefore do not measure an unaided language model.}
\end{table}

%% ============================================================
\section{Sweep Results: Noise Robustness (M3 vs.\ M4)}
\label{sec:noise}

\begin{table}[H]
\centering
\caption{$\Rsq$ statistics per noise level ($n=200$, 30 equations).
}
\label{tab:r2_noise}
\renewcommand{\arraystretch}{1.2}
\small
\begin{tabular}{l r r r r r r}
\toprule
& \multicolumn{3}{c}{\textcolor{m3blue}{\textbf{\EHD{} (M3)}}}
& \multicolumn{3}{c}{\textcolor{m4red}{\textbf{\HSL{} (M4)}}}\\
\cmidrule(lr){2-4}\cmidrule(lr){5-7}
$\sigma$ & Median & Min & Std & Median & Min & Std\\
\midrule
0\%    & 1.000000 & 0.999979 & $4.8\times10^{-6}$ & 1.0000000 & 1.000000 & $0.0$\\
0.05\% & 0.9999998 & 0.997348 & 0.0010 & 0.9999998 & 0.997807 & 0.0006\\
0.1\%  & 0.9999998 & 0.992748 & 0.0016 & 0.9999998 & 0.997810 & 0.0006\\
0.5\%  & 0.9999998 & 0.997504 & 0.0010 & 0.9999998 & 0.997789 & 0.0006\\
1\%    & 0.9999998 & 0.997325 & 0.0010 & 0.9999998 & 0.997720 & 0.0007\\
\bottomrule
\end{tabular}
\end{table}

\paragraph{Key observations.} The tables in this section (\ref{tab:r2_noise}, \ref{tab:rr_noise}, \ref{tab:time_noise}) are computed from the raw noise-sweep source files (\texttt{noise\_sweep\_20260812\_155519\_nshards03.json}, \texttt{...\_155520\_nshards05.json}, \texttt{...\_155532\_nshards01.json}, \texttt{...\_155541\_nshards02.json}, \texttt{...\_155544\_nshards04.json}), cross-checked against their raw per-test companions \texttt{protocol\_core\_noiseless\_20260812\_170344.json} and four \texttt{protocol\_core\_noisy\_*} files with matching timestamps to the second. The noise levels recorded in these files are $\sigma\in\{0, 0.05, 0.1, 0.5, 1\}\%$. M3 exhibits a std of $\sim$0.001--0.0016 across $\sigma>0$, higher
than M4's $\sim$0.0006--0.0007. M4 shows zero-variance convergence at
$\sigma=0\%$: all 30 equations reach $\Rsq=1.000000$ exactly, because the
LLM-first path generates the exact formula and the NN residual correction
drives the fit to machine precision. \textbf{M3 does not show the same
noiseless convergence}: at $\sigma=0\%$ M3's strict recovery rate
(threshold $\Rsq\geq0.999999$, per the source file's own recorded
\texttt{threshold\_used}) is \textbf{76.7\% (23/30)}, not 100\% --- see \S\ref{tab:rr_noise} below.

\begin{table}[H]
\centering
\caption{Recovery rate and catastrophic failure count per noise level ($n=200$).
Recovery uses each shard's recorded threshold: $\Rsq\geq0.999999$ at
$\sigma=0\%$, $\Rsq\geq0.995$ at $\sigma=0.05$, $0.1$ and $0.5\%$, and
$\Rsq\geq0.99$ at $\sigma=1\%$, so rates are not on a single scale across
rows.
}
\label{tab:rr_noise}
\small
\begin{tabular}{lrrrr}
\toprule
$\sigma$ & M3 Recovery & M3 Catastrophic & M4 Recovery & M4 Catastrophic\\
\midrule
0\%    & \textbf{76.7\%$^\ddagger$} & 0 & 100.0\% & 0\\
0.05\% & 100.0\%          & 0 & 100.0\% & 0\\
0.1\%  & 96.7\%           & 0 & 100.0\% & 0\\
0.5\%  & 100.0\%          & 0 & 100.0\% & 0\\
1\%    & 100.0\%          & 0 & 100.0\% & 0\\
\bottomrule
\multicolumn{5}{l}{\footnotesize$^\ddagger$ Source file \texttt{noise\_sweep\_..\_nshards01.json}, $\sigma=0$,}\\
\multicolumn{5}{l}{\footnotesize \texttt{threshold\_used}=0.999999: 23/30. M3 is \emph{not} perfect on noiseless recovery.}\\
\end{tabular}
\end{table}

\begin{table}[H]
\centering
\caption{Average per-equation computation time (seconds) per noise level.
Every cell is the per-equation \texttt{time} field averaged across all 30 tests in the five matching raw runner logs: \texttt{protocol\_core\_noiseless\_20260812\_170344.json} ($\sigma=0\%$) and \texttt{protocol\_core\_noisy\_20260812\_173425.json} ($\sigma=0.05\%$), \texttt{...\_172958.json} ($\sigma=0.1\%$), \texttt{...\_173523.json} ($\sigma=0.5\%$), and \texttt{...\_175318.json} ($\sigma=1\%$). The maximum M3 value across all 150 equation-runs (5 files $\times$ 30 equations) is 7.77\,s.}
\label{tab:time_noise}
\small
\begin{tabular}{lrrl}
\toprule
$\sigma$ & \textcolor{m3blue}{M3 avg (s)} & \textcolor{m4red}{M4 avg (s)} & Speedup\\
\midrule
0\%    & 7.55 & 6.92 & $1.1\times$\\
0.05\% & 7.66 & 6.55 & $1.2\times$\\
0.1\%  & 7.44 & 6.46 & $1.2\times$\\
0.5\%  & 7.77 & 6.56 & $1.2\times$\\
1\%    & 7.01 & 6.33 & $1.1\times$\\
\bottomrule
\end{tabular}
\end{table}

\paragraph{Note on M3 noiseless cost.} The source file for $\sigma=0\%$ records a mean M3 time of 7.55\,s, in line with every other noise level in \Cref{tab:time_noise}; none shows an exhaustive-search spike. The per-equation averages were re-derived from the raw runner logs (see the note on \Cref{tab:time_noise}) and agree with the table to two decimal places at all five noise levels. Across all
five real noise levels, M3 and M4 have comparable, stable per-equation
runtimes (M3: 7.0--7.8\,s; M4: 6.3--6.9\,s), with M4 consistently
$\sim$1.1--1.2$\times$ faster than M3.

%% ============================================================
\section{Sweep Results: Sample Complexity (M3 vs.\ M4)}
\label{sec:sc}

\begin{table}[H]
\centering
\caption{$\Rsq$ and RMSE per sample size (noiseless, 30 equations;
recovery thresholds used by the runner: $0.999$ at $n=50$, $0.9995$ at
$n=100$, $0.9999$ at $n\geq200$).
All six sample-size shards (\texttt{sample\_complexity\_20260812\_*\_nshards01--06.json}, $n=50$, $100$, $200$, $500$, $750$, $1000$, all \texttt{mode}: \texttt{noiseless}) were checked cell by cell against each shard's \texttt{per\_equation} data.}
\label{tab:sc_metrics}
\renewcommand{\arraystretch}{1.2}
\small
\begin{tabular}{r r r r r r r}
\toprule
& \multicolumn{3}{c}{\textcolor{m3blue}{\textbf{\EHD{} (M3)}}}
& \multicolumn{3}{c}{\textcolor{m4red}{\textbf{\HSL{} (M4)}}}\\
\cmidrule(lr){2-4}\cmidrule(lr){5-7}
$n$ & Med $\Rsq$ & Min $\Rsq$ & Med RMSE & Med $\Rsq$ & Min $\Rsq$ & Med RMSE\\
\midrule
  50 & 1.000000 & 0.9982 & 0.0006 & 1.0000000 & 1.0000 & 0.0000\\
 100 & 1.000000 & 0.9997 & 0.0011 & 1.0000000 & 1.0000 & 0.0000\\
 200 & 1.000000 & 1.0000 & 0.0009 & 1.0000000 & 1.0000 & 0.0000\\
 500 & 1.000000 & 0.9999 & 0.0016 & 1.0000000 & 1.0000 & 0.0000\\
 750 & 1.000000 & 0.9997 & 0.0021 & 1.0000000 & 1.0000 & 0.0000\\
1000 & 1.000000 & 0.9999 & 0.0013 & 1.0000000 & 1.0000 & 0.0000\\
\bottomrule
\end{tabular}
\end{table}

\HSL{} is exact (min $\Rsq=1.0000$, RMSE$=0.0000$) at every tested
sample size. \EHD{} reaches median $\Rsq=1.000000$ by $n=50$, but its
\emph{recovery rate} (fraction $\geq0.9999$) is 96.7\% (29/30) at
$n=50$, reaches 100\% at $n=100$ and $n=200$, then dips back to 96.7\%
at $n=500$, $750$, and $1000$ --- i.e.\ one equation stays just below
threshold at large $n$ rather than the metric monotonically improving
with sample size. ``Both converge at $n=50$'' is accurate for
\emph{median} $\Rsq$ but overstates \EHD's recovery-rate behaviour.

%% ============================================================
\section{Convergence Analysis}
\label{sec:convergence}

M3 shows stable convergence: $\Rsq$ std $\leq 0.001$ for $\sigma>0$.
M4 exhibits zero-variance convergence (std($\Rsq$) = 0.000) at all noise
levels in v2.

Both methods reach median $\Rsq=1.000000$ by $n=50$ in the noiseless sweep (all six sample-size
shards now verified; \Cref{tab:sc_metrics}).

\paragraph{Convergence boundary.}
The following conditions are needed for comprehensive generalisation claims:
(i)~higher-dimensional inputs ($\geq 6$ variables);
(ii)~PDE-form targets;
(iii)~extension to the full AI~Feynman dataset ($>100$ equations).

%% ============================================================
\section{Win Rate and Architecture Analysis}
\label{sec:winrate}

\begin{table}[H]
\centering
\caption{Head-to-head win rates (M3 vs.\ M4) on $\Rsq$. Noise sweep: 150 paired
comparisons (5 noise levels $\times$ 30 equations; pairs matched by shard and
test description). Sample-complexity (SC) sweep: 180 paired comparisons (6 sample
sizes $\times$ 30 equations; pairs matched by sample size and equation, from the
six \texttt{sample\_complexity\_*.csv} shards). Regenerated with
\texttt{compute\_winrate.py}. A pair is \emph{tied} if both systems have
$\Rsq>0.9999$ (the paper's original tie rule); otherwise the higher $\Rsq$ wins.
}
\label{tab:winrate}
\small
\begin{tabular}{l r r}
\toprule
\textbf{Outcome} & \textbf{Noise sweep} & \textbf{SC sweep}\\
\midrule
M3 strictly higher $\Rsq$ & 0/150 (0.0\%) & 0/180 (0.0\%)\\
M4 strictly higher $\Rsq$ & 28/150 (18.7\%) & 8/180 (4.4\%)\\
Tied (both $\Rsq>0.9999$)  & 122/150 (81.3\%) & 172/180 (95.6\%)\\
\bottomrule
\end{tabular}
\end{table}

\paragraph{Sensitivity to the tie rule.}
The split depends strongly on how ties are defined (noise sweep, $n=150$).
Under exact floating-point comparison M3 is nominally higher on 88 pairs and M4 on
62, but most of these differences are below $10^{-9}$ and reflect numerical noise
rather than a real accuracy gap. With an absolute tolerance of $10^{-6}$, M3 is
higher on 0, M4 on 35 (23.3\%) and 115 are tied (76.7\%); at $10^{-4}$ the figures
are 0 / 23 (15.3\%) / 127 (84.7\%); at $10^{-3}$ they are 0 / 19 (12.7\%) /
131 (87.3\%). Under no tolerance of $10^{-6}$ or looser does M3 exceed M4 on any
pair, and M4 is strictly higher on 12.7--23.3\% of pairs depending on the
tolerance. M4 therefore never loses to M3 beyond numerical noise on the noise
sweep, and most comparisons are ties.

The sample-complexity sweep ($n=180$) shows the same pattern. Under exact
comparison M3 is higher on 0 pairs, M4 on 179 and 1 is tied; with an absolute
tolerance of $10^{-6}$ the figures are 0 / 38 (21.1\%) / 142 (78.9\%); at
$10^{-4}$ they are 0 / 8 (4.4\%) / 172 (95.6\%); at $10^{-3}$ they are 0 / 1
(0.6\%) / 179 (99.4\%). M4's 8 strict wins under the $\Rsq>0.9999$ rule fall at
$n=50$ (2), $n=100$ (3), $n=200$ (0), $n=500$ (1), $n=750$ (1) and $n=1000$ (1).
M3 is never higher than M4 under any tie rule, including exact comparison.

\begin{table}[H]
\centering
\caption{Internal routing statistics --- evidence that M3 and M4 are genuinely
         distinct architectures. Regenerated by \texttt{generate\_tables.py}
         (\texttt{arch.tex}) from the five \texttt{protocol\_core} result
         files (one noiseless and four noisy runs,
         $\sigma\in\{0,0.05,0.1,0.5,1\}\%$; $5$ noise levels $\times$ $30$
         equations $= 150$ records per system). ``Ensemble-first'' (M3) and
         ``LLM-first'' (M4) denote the share of records whose recorded
         \texttt{metadata.decision} routed to that branch; this is a routing
         share, not a solve rate, and the \texttt{success} flag is not used.
         The pool mixes one noiseless run ($\sigma=0$) with four noisy runs.
         These counts supersede the earlier $92.7\%$/$74.7\%$ figures
         ($139/150$, $112/150$), which could not be reproduced from any
         available result file under any routing, \texttt{nn\_applied}, or
         $\Rsq$-threshold rule tried.}
\label{tab:arch}
\small
\begin{tabular}{L{3.5cm} r r}
\toprule
\textbf{Decision path} & \textbf{M3 (\EHD)} & \textbf{M4 (\HSL)}\\
\midrule
Primary path (routing share)
                & 140/150 (93.3\%) ensemble-first & 123/150 (82.0\%) LLM-first\\
Alternative path
                & 10/150 (6.7\%) NN              & 27/150 (18.0\%) ensemble\\
Secondary path  & NN refinement                  & NN residual correction\\
Architecture type & Ensemble-first               & LLM-first + NN-refinement\\
\bottomrule
\end{tabular}
\end{table}

\begin{table}[H]
\centering
\caption{Routing by noise level (30 equations per level), same source as
         \Cref{tab:arch} (\texttt{arch\_by\_noise.tex}). M3 routes identically at every level; M4 routes all
         noiseless records to the LLM but $20$--$27\%$ of noisy records to the
         ensemble.}
\label{tab:arch_by_noise}
\small
\begin{tabular}{l r r r r}
\toprule
$\sigma$ (\%) & M3 ensemble & M3 NN & M4 LLM & M4 ensemble\\
\midrule
0      & 28 & 2 & 30 & 0\\
0.05   & 28 & 2 & 22 & 8\\
0.1    & 28 & 2 & 24 & 6\\
0.5    & 28 & 2 & 23 & 7\\
1      & 28 & 2 & 24 & 6\\
\midrule
Total  & 140 & 10 & 123 & 27\\
\bottomrule
\end{tabular}
\end{table}

\paragraph{When to prefer M3 vs.\ M4.}
\begin{itemize}
\item \textbf{Prefer M3} when exact closed-form symbolic expressions are
  required, time is not a constraint, or noiseless exact symbolic recovery
  is the primary objective.
\item \textbf{Prefer M4} when wall-clock time is a factor
  (M4 is $\sim$1.1--1.2$\times$ faster than M3; \S\ref{sec:noise}), near-zero RMSE is prioritised, or $\sigma>0$ conditions are
  expected.
\end{itemize}

%% ============================================================
\section{Root-Cause Analysis: Newton's Gravity Measurement Bug}
\label{sec:bugfix}

\subsection{Symptom}

M4 reported $\Rsq \in [0.42, 0.87]$ on Newton's gravitational force across
all 11 tested conditions in v1, despite the LLM generating the correct formula
on every run.

\subsection{Root Cause}

\texttt{evaluate\_llm\_formula} used:
\begin{verbatim}
r2 = 1 - (ss_res / ss_tot) if ss_tot > 1e-10 else 0
\end{verbatim}
For Newton's gravity at tested parameter ranges, outputs are $\sim 10^{-11}$\,N,
giving $\mathrm{ss\_tot} \approx 1.8 \times 10^{-18}$ — eight orders of
magnitude below the threshold.  The function always returned
\texttt{r2 = 0.0}, triggering NN fallback.

\subsection{Fix Applied}

\begin{verbatim}
_tol = 1e-10 * float(np.max(np.abs(y_true))**2) * len(y_true)
r2   = 1 - (ss_res / ss_tot)
        if ss_tot > max(_tol, 1e-300) else 0
\end{verbatim}

The same fix was applied at three code locations:
\texttt{evaluate\_llm\_formula} (line 462),
\texttt{train\_nn} (line 384), and
the ensemble blend $\Rsq$ calculation (line 700).

\subsection{Additional Fixes in the Same Commit}

\begin{itemize}
  \item \texttt{force\_llm} flag: now correctly bypasses NN training.
  \item \texttt{PureLLMBaseline}: validates non-empty \texttt{python\_code}.
  \item Training epochs: corrected from 300 to 1000 (scheduler requires $\geq 1000$).
  \item Log-space guard: unconditionally enables log-space training when
    \texttt{y\_range > 3.0} decades.
\end{itemize}

\paragraph{General contribution.}
Any SR benchmark harness using absolute sum-of-squares thresholds will
silently misclassify correct formulas on equations with output values far
from unity.  The relative threshold fix should be adopted as standard practice.

%% ============================================================
\section{Discussion}
\label{suppb:discussion}

\subsection{Hardcoding and LLM Benchmark Integrity}

Of \PureLLM{}'s 30 noiseless rows, 11 (36.7\%) did not call the language model: 8 return the benchmark's own
ground-truth expression and 3 use a fixed form with parameters fitted to the data
(\Cref{tab:hardcoded}). In the 170344 file, \EHD{}'s returned formula is character-for-character identical to \PureLLM{}'s on all 11 of these bypassed rows, and on only 2 of the 19 live rows; the same holds in \texttt{benchmark\_results.json} (11 of 11 bypassed, 3 of 19 live). The cause is in the code: \EHD{} delegates formula generation to \texttt{PureLLMBaseline.generate\_formula}, which tries the hardcoded formulas before any language-model request; \EHD{}'s pass count on the bypassed rows (6 of 11, against 17 of 19 on live rows) also depends on its NN and ensemble logic. The all-domains LLM+NN hybrid delegates in the same way, but its stored output has no formula, so whether the bypass fired for it cannot be checked from the JSON. \PureLLM{} passes 29 of 30 rows (96.7\%) overall, 11 of 11 bypassed rows and 18 of 19
(94.7\%) \emph{live} rows, i.e.\ rows where the model was actually called. The live figure is not an unbiased
estimate of unaided performance: according to comments in the baseline source, the bypass branches were added for
equations on which the model had repeatedly failed, so excluding them favours \PureLLM{}. Counting every bypassed
row as a failure gives 18/30, so the unaided pass count lies between 18/30 and 29/30; the baseline has not been rerun
with the bypass disabled.
The recount over the per-domain files, \texttt{benchmark\_results.json} and \texttt{benchmark\_results\_pca\_4060.json} gives 29/30 and 18/19 on both splits. The \texttt{170344} file was not part of the recount; the one \PureLLM{} row below threshold in \texttt{benchmark\_results.json} is Curie's law, which returned no formula there ($\Rsq=0$, a failed or empty generation) and passes in the \texttt{170344} file. Which file is primary has not been decided.
Explicit hardcoding audits should be standard in any
benchmark involving LLM-based SR.

\subsection{M3 vs.\ M4 Trade-off}

With the measurement bug resolved, M3 and M4 are \textbf{close to
equivalent on $\Rsq$}: both reach median $\Rsq = 1.000000$ at every
tested noisy condition and sample size, though M3's \emph{recovery
rate} dips to 96.7\% at $\sigma=0.1\%$ (\Cref{tab:rr_noise}) and at
$n=500/750/1000$ (\Cref{tab:sc_metrics}), while M4 is 100\% throughout.
The distinction is: M3 is interpretability-optimised (exact symbolic
formulas); M4 is efficiency-optimised ($\sim$1.1--1.2$\times$ faster,
near-zero RMSE).

Under the $\Rsq>0.9999$ rule the tie rate is $81.3\%$ on the noise
sweep (122/150) and $95.6\%$ on the sample-complexity sweep (172/180), or
$89.1\%$ over the 330 pairs combined (\Cref{tab:winrate}). Direct computation
confirms that, on the noise sweep, M4 never loses to M3 on $\Rsq$ beyond
numerical noise ($|\Delta\Rsq|\le10^{-6}$), and that on the sample-complexity
sweep M3 is never higher than M4 under any tie rule.

%% ============================================================
\section{Limitations}
\label{suppb:limitations}

\begin{enumerate}
  \item \textbf{Benchmark scope.} 30 equations, 2--5 input variables.
    Higher-dimensional and PDE targets remain untested.
  \item \textbf{Memorisation scope.} The hardcoding audit detects verbatim
    recall but not partial memorisation.
  \item \textbf{Baseline breadth.} Only M3 and M4 covered by the sweep;
    adding PySR and DSR sweeps would strengthen generalisation claims.
  \item \textbf{M3 noiseless recovery.}  At $\sigma=0\%$, M3's strict recovery rate is 76.7\%
    (23/30, \Cref{tab:rr_noise}), not 100\%; this should be treated as an open finding
    requiring investigation, not merely a timing/framing issue.
  \item \textbf{No formula complexity metric.} Expression tree depth and
    operator count should accompany $\Rsq$ in future evaluations.
  \item \textbf{[Resolved, confirmed 2026-09-27.]} The project's CI
    dashboard (July~9 run) should not be cited: its own upstream notebooks
    (NB-01--NB-06) reported ``not found'' for that run, while its summary
    asserted ``0 open issues'' and specific claims resolved.
    \texttt{build\_report.py}'s \texttt{notebook\_execution\_status()}
    (commit \texttt{9d7a833f}, 2026-07-29, tagged
    \texttt{FIX-DASHBOARD-STALE-GREEN}) now checks per-cell
    \texttt{execution\_count} and hard-fails (\texttt{sys.exit(1)}) rather
    than publish a stale green summary over incomplete notebook output;
    confirmed present and active in the live repository. The July~9
    dashboard itself remains uncitable, but the pipeline defect is fixed
    and cannot recur silently.
  \item \textbf{[Resolved.]} The legacy exp2 baseline-lock threshold
    fix ($R^2 \geq 0.999999$ replacing $R^2 \geq 0.9999$, which had
    inflated 74/90 to 81/90) is now live-confirmed at 74/90 against a
    fresh CI run, and reproduced independently a second time via a
    patched Gate~C schema check; 81/90 should not be cited as current.
\end{enumerate}

%% ============================================================
\section{JMLR Readiness Assessment}
\label{sec:jmlr}

\noindent\textbf{Note (v3.0 update, July 2026).}
The readiness metrics in Table~\ref{tab:jmlr} below are based on the v2
benchmark run (30 benchmark equations, noiseless/noise sweep protocol).
The main paper (\texttt{jmlr\_paper\_main.tex}) now reports v3.0
results on 74 DeFi tasks under a 40\,\%/60\,\% aggressive extrapolation
split (the result file does not record whether the single-axis or PCA-directed variant was used): HypatiaX's near-perfect-success figures are \textbf{90.5\,\%} (67/74) raw
uncorrected / \textbf{59.5\,\%} (44/74) corrected --- resolved by
cross-reference to the main paper's Hybrid Decision-Attribution Bug
section (\texttt{jmlr\_paper\_main.tex}, \S~``The Hybrid
Decision-Attribution Bug''), which shows that the corrected version does
\emph{not} eliminate the 5 Pure-LLM catastrophic failures: the hybrid
routes to the LLM on all 5 and its zero-catastrophic report was masked by the
same attribution bug (corrected catastrophic count: 5/74, identical to Pure LLM;
main paper Table~\texttt{tab:main\_results} recomputed 2026-09-28 from the same file). \textbf{Independently
re-verified (2026-09-27)} by direct recomputation against the canonical
seed-42 file retrieved from
\texttt{github.com/sednabcn/LLM-HypatiaX-DEV} (md5
\texttt{cbf503d0b3985b4c1bccf7ed18e0316b}, canonical \texttt{hypatiax/data}; recomputed 2026-09-29); see also Post-Submission
Report, \S~``Item 2.'' These are a \emph{different and more stringent}
evaluation than the noiseless recovery metrics here; the v2 benchmark readiness
assessment remains valid for this document's scope.

\begin{table}[ht]
\centering
\caption{JMLR readiness assessment (v2 benchmark results, post bug-fix).
  See note above for v3.0 DeFi results reported in the main paper.}
\label{tab:jmlr}
\small
\renewcommand{\arraystretch}{1.3}
\begin{tabular}{l C{1.5cm} C{1.5cm} L{5.0cm}}
\toprule
\textbf{Criterion} & \textbf{M3} & \textbf{M4} & \textbf{Required action}\\
\midrule
Accuracy ($\Rsq$)       & \textcolor{passgreen}{\textbf{PASS}} & \textcolor{passgreen}{\textbf{PASS}} &
  Both: median $\Rsq=1.000000$\\
Noise robustness        & \textcolor{warnAmber}{\textbf{COND.}} & \textcolor{passgreen}{\textbf{STRONG}} &
   M3 76.7\% at $\sigma{=}0$, 96.7--100\% at $\sigma>0$ (\Cref{tab:rr_noise}); M4 100\% at all $\sigma$\\
Data efficiency         & \textcolor{passgreen}{\textbf{STRONG}} & \textcolor{passgreen}{\textbf{STRONG}} &
  Convergence at $n=50$ for both (noiseless only)\\
Catastrophic failures   & \textcolor{passgreen}{\textbf{CLEAN}} & \textcolor{passgreen}{\textbf{CLEAN}} &
  0/30 for both; document Newton bug\\
Computational cost      & \textcolor{passgreen}{\textbf{STRONG}} & \textcolor{passgreen}{\textbf{STRONG}} &
   Both $<8$\,s/eq across all $\sigma$ (\Cref{tab:time_noise}); no evidence of an offline-search spike\\
Statistical significance & \textcolor{passgreen}{\textbf{PASS}} & \textcolor{passgreen}{\textbf{PASS}} &
  Add Wilcoxon test; tie rates recomputed in \S\ref{sec:winrate} (81.3\% noise, 95.6\% SC)\\
Hardcoding disclosure   & \multicolumn{2}{c}{\textcolor{warnAmber}{\textbf{REQUIRED}}} &
  \PureLLM: disclose 11/30 hardcoded\\
Process/CI credibility  & \multicolumn{2}{c}{\textcolor{warnAmber}{\textbf{OPEN}}} &
  Baseline-lock now live-confirmed (74/90); do not cite July~9 dashboard until re-run (\S\ref{suppb:limitations})\\
\midrule
\textbf{Overall verdict} & \textcolor{passgreen}{\textbf{READY}} & \textcolor{passgreen}{\textbf{READY}} &
  Ready on accuracy/robustness grounds; one process item above remains open\\
\bottomrule
\end{tabular}
\end{table}

%% ============================================================
\section{Conclusions}
\label{sec:conclusions}

\begin{enumerate}
  \item \textbf{Both M3 and M4 achieve 100\% recovery} at all tested noisy
    conditions ($0.05$--$1\%$ noise) and all
    sample sizes ($50$--$1000$).
  \item \textbf{M4 reaches 100\% at $\sigma=0$} (noiseless, v2-corrected);
    \textbf{M3 reaches only 76.7\% (23/30) at $\sigma=0$}
    (\Cref{tab:rr_noise}). \HDS{} (M6, tools track) trails at 70.0\% (21/30)
    under the separate six-method noiseless protocol (30 equations;
    \Cref{tab:overall}), with nine failing equations (Arrhenius rate
    constant, Nernst equation, dielectric polarisation, capacitor energy,
    Coulomb's law, kinetic energy, Gaussian/normal PDF, Bose--Einstein
    occupation number, Fermi--Dirac occupation number).

  \item \textbf{M4 is $\sim$1.1--1.2$\times$ faster} with near-zero RMSE
    ; the throughput advantage is modest, not
    decisive.
  \item \textbf{\EHD{} sits close to AI~Feynman~2.0's published figure}:
    76.7\% vs.\ 79.3\% ($-2.6$\,pp) on this work's own custom 30-equation
    benchmark. \textbf{\HDS{} falls further below it}: 70.0\% ($-9.3$\,pp).
      Published baselines were evaluated on all 100
    Feynman equations under SRBench, a different equation set entirely;
    the AI Feynman figure is shown for general context only, not as a
    head-to-head comparison.
  \item \textbf{The Newton failure in v1 was a measurement bug}: absolute
    $\Rsq$ threshold misclassified correct formulas for tiny-scale equations.
    The relative threshold fix is a general contribution to SR benchmarking.
  \item \textbf{Four methodological issues corrected}: undisclosed hardcoding,
    scale-biased \rmse, missing significance tests, mismatched thresholds.
\end{enumerate}

%% ── Bibliography ─────────────────────────────────────────────────────────────
\bibliography{references}

%% ============================================================
\appendix

\section{Reproducibility Figure Generation Inventory}
\label{app:figures}

\noindent This appendix documents the auxiliary figures produced by
\texttt{generate\_figures.py} for reproducibility purposes. It is
distinct from the figures actually embedded in the compiled
manuscript (see \texttt{figures\_inventory.tex}, 5 files, listed in
\texttt{jmlr\_paper\_main.tex}).

\begin{table}[H]
\centering
\caption{Complete figure inventory; all files in \texttt{figures/}.
  All 13 generated figures must be regenerated from v2 result JSONs.}
\label{tab:figlist}
\small
\renewcommand{\arraystretch}{1.25}
\begin{tabular}{L{4.8cm} L{1.5cm} L{1.8cm} L{3.2cm}}
\toprule
\textbf{Filename} & \textbf{Sweep} & \textbf{Section} & \textbf{Content}\\
\midrule
\texttt{fig1\_r2\_vs\_noise.pdf}       & Noise & \ref{sec:noise}        & Median $\Rsq$ vs $\sigma$\\
\texttt{fig2\_rmse\_vs\_noise.pdf}     & Noise & \ref{sec:noise}        & Median RMSE vs $\sigma$\\
\texttt{fig3\_time\_vs\_noise.pdf}     & Noise & \ref{sec:noise}        & Avg time vs $\sigma$\\
\texttt{fig4\_r2\_vs\_n.pdf}           & SC    & \ref{sec:sc}           & Median $\Rsq$ vs $n$\\
\texttt{fig5\_rmse\_vs\_n.pdf}         & SC    & \ref{sec:sc}           & Median RMSE vs $n$\\
\texttt{fig6\_time\_vs\_n.pdf}         & SC    & \ref{sec:sc}           & Median time vs $n$\\
\texttt{fig7\_recovery\_vs\_noise.pdf} & Noise & \ref{sec:noise}        & Recovery rate vs $\sigma$\\
\texttt{fig8\_recovery\_vs\_n.pdf}     & SC    & \ref{sec:sc}           & Recovery rate vs $n$\\
\texttt{fig9\_minr2\_vs\_noise.pdf}    & Noise & \ref{sec:noise}        & Min $\Rsq$ vs $\sigma$\\
\texttt{fig10\_r2\_boxplot\_noise.pdf} & Noise & \ref{sec:convergence}  & Per-eq $\Rsq$ box plots\\
\texttt{fig11\_recovery\_heatmap.pdf}  & Both  & \ref{sec:convergence}  & Recovery heatmap $\sigma \times n$\\
\midrule
\texttt{fig\_runtime\_comparison.png}  & ---   & \ref{sec:noiseless}    & Runtime bar chart, 6 methods\\
\texttt{fig\_comparative\_table.png}   & ---   & \ref{sec:noiseless}    & Domain $\times$ method table\\
\bottomrule
\end{tabular}
\end{table}

\noindent\textbf{Generation note:} Regenerate from
\texttt{noise\_sweep\_20260316\_192711.json} and
\texttt{sample\_complexity\_20260316\_193447.json} (v2 corrected
harness; see Section~\ref{sec:r2_measurement_bugfix} for v1 vs.\ v2 provenance).

\noindent\textbf{Relationship to the submitted PDF:} the 13 files
above are reproducibility documentation only and are not
\verb|\includegraphics|'d anywhere in this document or in the main
manuscript. The compiled manuscript embeds a separate, disjoint set
of 5 figures directly via \verb|\includegraphics| (see
\texttt{jmlr\_paper\_main.tex}). Only those 5 files are bundled in
\texttt{submission/sources/figures/}; the 13 listed here should be
regenerated from source by anyone who wants them, not shipped in the
submission package.

\section{Method Abbreviations}
\label{app:abbrev}

\begin{tabular}{L{1.8cm} L{2.5cm} L{6.0cm}}
\toprule
\textbf{Short} & \textbf{Code} & \textbf{Full name}\\
\midrule
\PureLLM & M1 & PureLLM Baseline (core)\\
\INN     & M2 & ImprovedNN (core)\\
\EHD     & M3 & EnhancedHybridSystemDeFi (core)\\
\HSL     & M4 & HybridSystemLLMNN all-domains (core)\\
\SEL     & M5 & SymbolicEngineWithLLM (tools)\\
\HDS     & M6 & HybridDiscoverySystem v50\_2 (tools)\\
\bottomrule
\end{tabular}

\section{PureLLM Required Disclosure}
\label{app:disclosure}

Required in any JMLR submission including \PureLLM{} results:

\begin{quote}
\textit{The PureLLM Baseline does not call the language model for 11 of 30
benchmark equations (\texttt{is\_hardcoded: true}): for 8 it returns the benchmark's own
ground-truth formula and for 3 a fixed form whose parameters are fitted to the data by
least squares. Its pass rate is 29/30 over all rows and 18/19 on the 19 live rows
(where the model was called). The live figure favours PureLLM because the bypass
was added for equations on which the model had failed, so the unaided pass count
lies between 18/30 and 29/30; the baseline has not been rerun with the bypass disabled.}

\end{quote}

% ====================================================================
% NEW APPENDICES (extracted from jmlr_paper_main.tex)
% ====================================================================
\section{Benchmark Reproducibility Statement}
\label{app:repro_statement}

\subsection*{Code, Data, and Protocol}

All experimental results in the main paper's \emph{Benchmark Design} section were produced
using \texttt{experiment\_protocol\_benchmark\_v2.py} (version~2.0,
2026-03-02), which will be released alongside the final paper.
The protocol file encodes all data-generation functions, random seeds, noise
levels, and recovery thresholds as class constants and is self-documenting:

\begin{verbatim}
python experiment_protocol_benchmark_v2.py --describe
\end{verbatim}

\noindent The 30-equation benchmark is defined in
\texttt{\_build\_feynman\_equations()} (lines 480--953 of the protocol file; the
function name uses the label ``feynman'', but the set is custom, not a Feynman subset).
Each equation specifies its identifier, ground-truth formula, variable
ranges, and generating function; data are produced analytically, eliminating
external dataset dependencies for the physics equations.

\subsection*{Exact Reproduction Commands}

\paragraph{Phase~2 (noisy, $R^2>0.995$).\textbf{[EDITORIAL NOTE: no table in this
document reports a noisy-protocol, six-method, 30-equation aggregate at the
0.995 threshold --- the original \texttt{tab:feynman30} label does not exist
anywhere in the paper or supplements. Please confirm whether this should
point to a specific table (e.g.\ a per-equation noisy breakdown analogous to
Table~\ref{tab:overall}) or be removed.]}}
\begin{verbatim}
python run_comparative_suite_benchmark_v2.py \
    --methods 1 --samples 200 --no-llm-cache
\end{verbatim}
Output: \texttt{protocol\_core\_noisy\_<timestamp>.json}.
The \texttt{--no-llm-cache} flag forces independent fresh API calls for all
30 Pure~LLM evaluations; this was verified via cache-clearing.

\paragraph{Phase~3 (noiseless, $R^2\geq0.999999$, Table~\ref{tab:overall}).}
The command and output below come from the one verified noiseless source file, \texttt{protocol\_core\_noiseless\_20260812\_170344.json}, whose internal \texttt{script} field is quoted directly.
\begin{verbatim}
python run_protocol_benchmark_core.py \
    --noiseless --threshold 0.999999     \
    --nn-seeds 3 --samples 200           \
    --method-timeout 900                 \
    --pysr-timeout 1100
\end{verbatim}
Script version: \texttt{run\_protocol\_benchmark\_core.py v2.2 (sign-fix +
log-widen + domain-guard + formula-hash + complexity-score + json-export +
extrap\_r2\_far-fix)}, per the source file's own \texttt{script} field.\\
Output: \texttt{protocol\_core\_noiseless\_20260812\_170344.json}
(timestamp 2026-08-12T17:03:44; 30 equations across 11 non-Feynman domains
--- biology, chemistry, electrochemistry, electromagnetism, electrostatics,
magnetism, mechanics, optics, probability, quantum, thermodynamics --- run
against all 6 methods, with per-method pass counts as read from that file, which
\Cref{tab:overall} reproduces except for \PureLLM{} (29/30 there, from the recount) (\PureLLM{} 30/30, \HSL{} 30/30, \SEL{} 28/30, \EHD{}
23/30, \HDS{} 21/30, \INN{} 0/30); wall-clock and hardware not recorded in
the source file, so no timing claim is made here.

\textbf{Note on timeout upgrade:} \texttt{--method-timeout} was set to 300\,s
for tests 1--18 and upgraded to 900\,s for tests 19--30. Arrhenius (test~4)
ran under the 300\,s limit and the harness recorded a 27574\,s wall-clock
before the run ended; this is recorded as a genuine failure independent of
the v2 bug fix. \textbf{Revised mechanism (Phase~3 audit, item~10b):} PySR is
invoked via \texttt{juliacall}, which embeds Julia in-process rather than
spawning a separate OS process, so there was no external Julia PID being
``killed.'' The 27574\,s figure instead reflects a gap in the harness's own
timeout wiring: the outer \texttt{ThreadPoolExecutor} timeout fires at
300\,s but abandons the background thread without a reference back to the
underlying call, which then continues to run unobserved. \textbf{Update
(Phase~3 audit, Item~10b, open):} The two descriptions above do not match: the root-cause paragraph states that PySR runs in-process via \texttt{juliacall} with no separate OS process, whereas the fix threads a live \texttt{Popen} handle to \texttt{\_kill\_process\_group}, which presupposes a separate process. Either the root-cause description is incomplete, the fix addresses a different, subprocess-based hang, or \texttt{verify\_proc\_box\_fix.py} exercises a different code path. Resolving this requires checking the \texttt{\_ProcBox} implementation and the verification script's PySR invocation against \texttt{run\_comparative\_suite\_benchmark\_v2\_FIXED.py}. The originally
stated fix (a thread-safe \texttt{\_ProcBox}
handle threading the live \texttt{Popen} reference for the in-flight PySR
subprocess up to the outer timeout handler, which then terminates it via
\texttt{\_kill\_process\_group}) is implemented in
\texttt{run\_comparative\_suite\_benchmark\_v2\_FIXED.py} and has been
wall-clock verified twice: once in a sandboxed reproduction using a
synthetic hang-prone subprocess, and once on the real Julia/PySR stack
(\texttt{verify\_proc\_box\_fix.py}, unmodified, run on the user's own
machine). Both verifications show the outer timeout now bounds real
wall-clock time rather than only the abandoned thread's timeout --- but
see the open question above regarding what mechanism this
actually verifies. See the
companion document \emph{Phase 3 Reproducibility Audit: Consolidated
Findings and Bug Recovery History} (Item~10b) for the
full diagnosis, including a superseded intermediate root-cause model.
Nernst (test~6) also timed out under the 300\,s default; counted as failure.

\subsection*{Random Seeds}

\begin{itemize}
  \item Base seed: \texttt{42} (class constant
    \texttt{BenchmarkProtocol.seed = 42}, line~1091 of the protocol file).
    Matches the Core~15 benchmark seed for cross-campaign consistency.
  \item Per-equation seed: \texttt{seed + offset $\times$ 7}, where
    \texttt{offset} is the 0-based index in the equation list
    (line~1208: \texttt{eq.generate(n, noise, seed=s + offset * 7)}).
    This ensures each of the 30 equations receives a distinct, reproducible
    seed regardless of call order.
  \item Neural network seeds (Phase~3): 3 independent seeds per equation;
    median $R^2$ reported. Exact seed values are logged in the output JSON.
\end{itemize}

\subsection*{Software Environment}

\begin{table}[H]
\centering
\caption{Software and hardware environment for benchmark runs
         (verified 18 March 2026 via \texttt{hardware\_info.py} and
         \texttt{pip show}).}
\label{tab:software_env}
\small
\begin{tabular}{ll}
\toprule
\textbf{Component} & \textbf{Version / details} \\
\midrule
Python           & 3.12.2 \\
PySR             & 1.5.9 \\
Julia runtime    & 1.12.2 \\
juliacall        & 0.9.26 (Python$\leftrightarrow$Julia bridge used by PySR) \\
NumPy            & 2.3.5 \\
SciPy            & 1.16.3 \\
SymPy            & 1.14.0 \\
Pandas           & 2.3.3 \\
Matplotlib       & 3.10.7 \\
Seaborn          & 0.13.2 \\
scikit-learn     & 1.7.2 \\
Anthropic SDK    & 0.73.0 (\texttt{claude-sonnet-5},
                   temperature~0.25) \\
\midrule
Hardware         & Intel Celeron T3100 @ 1.90\,GHz (2~cores);
                   3\,GB DDR2; no GPU \\
OS               & Kali GNU/Linux Rolling \\
PySR populations & 2 (hardware limit: 2-core CPU) \\
\bottomrule
\end{tabular}
\end{table}

\noindent\textbf{Note on Anthropic SDK version compatibility.}
\texttt{anthropic==0.73.0} is substantially newer than early-2024 releases.
The Messages API interface is stable across 0.18$+$ versions; earlier
installations may require minor call-site adjustments but the prompts and
response parsing are otherwise unchanged.

\noindent PySR uses Julia under the hood; the 1100\,s
\texttt{--pysr-timeout} (\texttt{FEYNMAN\_TIMEOUT}, the inner PySR fit
budget---distinct from the 900\,s outer \texttt{--method-timeout} above,
which governs the per-method wall-clock allowance, not the PySR fit itself)
was required to accommodate Julia JIT initialisation on the first run.
Subsequent equations in the same process are faster. The 27574\,s Arrhenius
figure is a harness-timeout artifact, not evidence that PySR's complexity
search fails to converge on single-variable flat-landscape equations; see
the revised mechanism note above.

\subsection*{Measurement Correction (v1\,$\to$\,v2)}
\label{sec:r2_measurement_bugfix}

A bug was identified on 2026-03-16 in \texttt{evaluate\_llm\_formula}
and corrected before final submission. The function used a hardcoded
absolute sum-of-squares guard (\texttt{ss\_tot > 1e-10}) that silently
returned $R^2=0$ for equations whose output magnitudes are far below unity
(Newton gravity $\sim\!10^{-11}$\,N, Zeeman splittings $\sim\!10^{-23}$\,J,
Coulomb forces at $\mu$C scale). The fix scales the threshold relative
to the data:

\begin{verbatim}
# v1 (buggy):
r2 = 1 - (ss_res / ss_tot) if ss_tot > 1e-10 else 0

# v2 (fixed):
_tol = 1e-10 * float(np.max(np.abs(y_true))**2) * len(y_true)
r2 = 1 - (ss_res / ss_tot) if ss_tot > max(_tol, 1e-300) else 0
\end{verbatim}

The same fix was applied at three locations:
\begin{enumerate}
  \item \texttt{evaluate\_llm\_formula} (line 462) — primary failure site
  \item \texttt{train\_nn} (line 384) — NN self-evaluation
  \item Ensemble-blend $R^2$ calculation (line 700) — candidate scoring
\end{enumerate}

Additional fixes in the same commit:
\begin{itemize}
  \item \texttt{force\_llm} flag: \texttt{hybrid\_predict} now correctly
    bypasses NN training when the LLM formula passes the threshold.
  \item \texttt{PureLLMBaseline} passthrough: \texttt{generate\_llm\_formula}
    now validates non-empty \texttt{python\_code} before accepting.
  \item Training epochs: corrected from 300 to 1000 to satisfy the learning
    rate scheduler's warm-restart requirement ($\geq\!1000$ epochs).
  \item Log-space guard: \texttt{\_is\_pole} heuristic now unconditionally
    enables log-space training when \texttt{y\_range > 3.0} decades,
    fixing suppression on wide-range power-law equations.
\end{itemize}

The \emph{Summary of Changes} table in Supplementary~A (\emph{Routing Improvements for Hybrid LLM--Neural-Network Systems}), and Section~\ref{sec:r2_measurement_bugfix} of this document, summarise all
aggregate recovery changes. All results reported in the paper use the
v2 corrected harness. The v1 JSON output files are retained for audit
purposes and are available upon request.

\subsection*{Per-Equation v2 Data}

The complete per-equation $R^2$ values under the v2 corrected harness are
stored in:
\begin{itemize}
  \item \texttt{protocol\_core\_noiseless\_20260304\_154510.json} —
    Phase~3 noiseless, all 6 methods, 30 equations.
  \item \texttt{noise\_sweep\_20260316\_192711.json} —
    M3/M4 noise sweep ($\sigma \in \{0, 0.05, 0.1, 0.5, 1\}\%$,
    , v2 corrected).
  \item \texttt{sample\_complexity\_20260316\_193447.json} —
    M3/M4 sample complexity sweep ($n \in \{50, 100, 200, 500, 750, 1000\}$,
    v2 corrected).
\end{itemize}
These files will be released as supplementary material. The v1 superseded
files (\texttt{noise\_sweep\_20260315\_091018.json},
\texttt{sample\_complexity\_20260315\_124310.json}) should not be used
for the final submission figures.

\subsection*{PureLLM Disclosure}

Following~\citet{memorization2024audit}, each PureLLM noiseless row was audited for whether the
language model was called. \textbf{11 of 30 PureLLM rows bypass the model
(\texttt{is\_hardcoded=True})}: Michaelis--Menten, Logistic growth,
Allometric scaling, Nernst, Clausius--Mossotti, Lorentz force, Gaussian PDF,
Zeeman energy, Bose--Einstein, Planck radiance, and Fourier heat. Eight of these
return the benchmark's ground-truth formula; Michaelis--Menten, Logistic growth and Allometric scaling use a fixed form with
least-squares parameters fitted to the data. PureLLM passes 29/30 over all rows and 18/19 on the 19 live rows
(where the model was called); because the bypass was added for equations on which the model had failed,
the live figure favours PureLLM and the unaided pass count lies between 18/30 and 29/30.

PureLLM is listed in all tables, and the main paper reports its comparisons with and without the bypassed rows.

A controlled re-run with \texttt{--no-llm-cache} confirmed that all
30 equations issued independent API calls and produced identical scores
(median $R^2=0.9976$, std $=2.5\times10^{-4}$), ruling out any
caching artefact.

\subsection*{Citation Key Corrections}

The following citation key conflicts have been resolved in \texttt{references.bib}
(updated 2026-06):

\begin{enumerate}
  \item \texttt{cranmer2023interpretable} renamed to \texttt{cranmer2023pysr}
        (canonical key matching all \texttt{\textbackslash cite} calls in the main paper;
        same paper arXiv:2305.01582). \textbf{Resolved.}
  \item \texttt{udrescu2020aifeynman} duplicate removed; \texttt{udrescu2020ai}
        retained as the single canonical key in \texttt{references.bib}.
        This supplementary document uses \texttt{udrescu2020aifeynman} in its own
        inline \texttt{thebibliography} and remains self-consistent. \textbf{Resolved.}
  \item \texttt{meidani2023snip} (arXiv preprint key) replaced by
        \texttt{meidani2024snip} (published ICLR~2024 key) in
        \texttt{references.bib}. \textbf{Resolved.}
  \item \texttt{lacava2021contemporary} entry confirmed present in
        \texttt{references.bib}. \textbf{No action needed.}
\end{enumerate}

\subsection*{SNIP Comparison Note}

SNIP~\citep{meidani2024snip} does \emph{not} report recovery rates.
It evaluates symbolic regression using a Pareto front of median
$R^2$ vs.\ equation complexity on 119 Feynman equations, reporting
median $R^2 = 0.882$ (vs.\ AI~Feynman baseline: $0.798$).
Placing SNIP's median $R^2$ alongside HypatiaX's recovery rate
(\% at $R^2>0.9999$) would be a category error.
The only valid comparison uses the shared median $R^2$ axis:
HypatiaX Hybrid DeFi achieves median $R^2 = 1.0000$ on our 30-equation
subset vs.\ SNIP's $0.882$ on the full 119-equation SRBench set
(different subsets; comparison is indicative only).
The Nguyen-12 benchmark table in the main paper provides this comparison.

% ============================================================================
% APPENDIX A: COMPLETE PROOF OF THEOREM
% ============================================================================


\section{Hyperparameter Settings}
\label{app:hyperparameters}

\begin{table}[H]
\centering
\caption{PySR Configuration for Symbolic Regression}
\label{tab:hyperparameters}
\begin{tabular}{@{}lc@{}}
\toprule
\textbf{Parameter} & \textbf{Value} \\
\midrule
Population size & 100 \\
Generations & 500 \\
Populations (parallel) & 2 (hardware limit: 2-core CPU) \\
Crossover probability & 0.8 \\
Mutation probability & 0.1 \\
Complexity penalty & 0.01 \\
Tournament size & 3 \\
Elitism & 5 \\
Max expression depth & 15 \\
Optimization iterations & 50 \\
\midrule
\textbf{Validation Thresholds} & \\
\midrule
Acceptance $R^2$ (Criterion 1) & 0.99 \\
Validation score (Criterion 1) & 30 \\
Acceptance $R^2$ (Criterion 2) & 0.95 \\
Validation score (Criterion 2) & 80 \\
Extrapolation $R^2$ & 0.80 \\
Numerical stability ($\alpha$) & 0.95 \\
\bottomrule
\end{tabular}
\end{table}

\begin{table}[H]
\centering
\caption{Complete Hyperparameter Specifications for All Methods}
\label{tab:hyperparameters_complete}
\small
\begin{tabular}{lll}
\toprule
\textbf{Method} & \textbf{Parameter} & \textbf{Value} \\
\midrule
\multirow{8}{*}{Neural Network}
& Architecture & [64, 32, 1] \\
& Activation & ReLU \\
& Optimizer & Adam \\
& Learning rate & 0.01 \\
& Batch size & Full batch (160 samples) \\
& Epochs & 200 \\
& Weight initialization & Xavier uniform \\
& Dropout & None \\
\midrule
\multirow{7}{*}{Pure LLM}
& Model & claude-sonnet-5 \textbf{[corrected: was
  claude-sonnet-4-20250514]} \\
& Temperature & 0.25 \textbf{[corrected: was 0.0 (deterministic)]} \\
& Max tokens & 1024 \\
& Top-p & 1.0 \\
& Few-shot examples & 2-3 per domain \\
& Prompt template & Domain-specific \\
& Retry on failure & 3 attempts \\
\midrule
\multirow{12}{*}{Hybrid v50\_2}
& SR iterations & 50 \\
& SR populations & 2 (hardware limit: 2-core CPU) \\
& Binary operators & [+, -, *, /] \\
& Unary operators & [exp, log, sqrt, sin, cos] \\
& Parsimony coefficient & 0.01 \\
& Tournament size & 10 \\
& Mutation rate & 0.1 \\
& Crossover rate & 0.9 \\
& LLM candidates & 5 \\
& LLM temperature & 0.25 \textbf{[corrected: was 0.3]} \\
& Max complexity & 30 \\
& Validation threshold ($R^2$) & 0.99 \\
\bottomrule
\end{tabular}
\end{table}

% ============================================================================
% APPENDIX C: DETAILED EXPERIMENTAL RESULTS
% ============================================================================


\section{Detailed Experimental Results}
\label{app:detailed_results}

\subsection{Pure LLM Baseline: Per-Case Results}

\input{tables/llm_detailed.tex}

\subsection{DeFi-Optimized HypatiaX: Per-Case Results}

\input{tables/defi_detailed.tex}

\subsection{Complete Test Suite Details}

\begin{table}[H]
\centering
\caption{Ground Truth Formulas with Exact Constants}
\label{tab:formulas_detailed}
\scriptsize
\begin{tabular}{lll}
\toprule
\textbf{Equation} & \textbf{Formula} & \textbf{Constants} \\
\midrule
Arrhenius & $k = A\exp(-E_a/(RT))$ & $A=10^{11}$, $E_a=80000$ J/mol, $R=8.314$ J/(mol$\cdot$K) \\
Henderson-Hasselbalch & $\text{pH} = \text{p}K_a + \log_{10}([A^-]/[\text{HA}])$ & $\text{p}K_a = 6.5$ \\
Rate Law & $r = k[A]^2[B]$ & $k = 0.5$ M$^{-2}$s$^{-1}$ \\
Allometric Scaling & $Y = aM^b$ & $a = 3.5$, $b = 0.75$ \\
Michaelis-Menten & $v = V_{\max}[S]/(K_m + [S])$ & $V_{\max} = 50$ $\mu$M/s, $K_m = 10$ $\mu$M \\
Logistic Growth & $dN/dt = rN(1-N/K)$ & $r = 0.3$ day$^{-1}$, $K = 1000$ \\
Kinetic Energy & $E = \frac{1}{2}mv^2$ & Dimensionless \\
Gravitational Force & $F = Gm_1m_2/r^2$ & $G = 6.674 \times 10^{-11}$ N$\cdot$m$^2$/kg$^2$ \\
Ideal Gas Law & $P = nRT/V$ & $R = 8.314$ J/(mol$\cdot$K) \\
Impermanent Loss & $\text{IL} = 2\sqrt{r}/(1+r) - 1$ & Dimensionless \\
Price Impact & $\Delta p = \Delta x/(x + \Delta x)$ & Dimensionless \\
Constant Product & $y = k/x$ & $k = 10^6$ \\
VaR 95\% & $\text{VaR} = P\sigma \cdot z_{0.95}$ & $z_{0.95} = 1.645$ \\
Liquidation Price & $p_{\text{liq}} = p_0(1 - 1/(L \cdot m))$ & $m = 0.8$ (maintenance margin) \\
Portfolio Variance & $\sigma_p = \sqrt{\sigma_1^2 + \sigma_2^2 + 2\rho\sigma_1\sigma_2}$ & Dimensionless \\
\bottomrule
\end{tabular}
\end{table}

% ============================================================================
% APPENDIX D: STATISTICAL TEST DETAILS
% ============================================================================

\section{Statistical Test Details}
\label{app:statistical_tests}

\subsection*{Experiment 1 (exp1\_five, Core-15)}

\subsection{Mann-Whitney U Test for Extrapolation Performance}

For medium extrapolation regime (2$\times$ training range):

\paragraph{Test Setup:}
\begin{itemize}
\item \textbf{Sample sizes}: $n_{\text{Hybrid v50\_2}} = 12$, $n_{\text{Neural Network}} = 15$
\item \textbf{Null hypothesis}: $H_0: E_{\text{Hybrid v50\_2}} \geq E_{\text{Neural Network}}$ (Hybrid v50\_2 is not better)
\item \textbf{Alternative hypothesis}: $H_1: E_{\text{Hybrid v50\_2}} < E_{\text{Neural Network}}$ (Hybrid v50\_2 is better)
\item \textbf{Significance level}: $\alpha = 0.05$ (one-tailed test)
\end{itemize}

\paragraph{Descriptive Statistics:}
\begin{itemize}
\item \textbf{Hybrid v50\_2}: Mean = 1.362e+26\%, Median = 2415.2\%, SD = 4.717e+26\%, Range = [1.3, 1.634e+27]
\item \textbf{Neural Network}: Mean = 8899923.4\%, Median = 13283.9\%, SD = 33577002.8\%, Range = [216.4, 130261180.6]
\end{itemize}

\paragraph{Test Results:}
\begin{itemize}
\item \textbf{Test statistic}: $U = 74.0$
\item \textbf{P-value}: $p = 2.25e-01$ (one-tailed, normal approximation with tie correction -- see \texttt{\_render\_stat\_test\_tex()} in \texttt{generate\_tables.py})
\item \textbf{Conclusion}: $p \geq 0.05 \Rightarrow$ Fail to reject $H_0$ at $\alpha = 0.05$
\end{itemize}

\subsection{Effect Size Calculation (Cohen's $d$)}

\begin{align}
\mu_{\text{Neural Network}} &= 8899923.4\%, \quad \sigma_{\text{Neural Network}} = 33577002.8\% \\
\mu_{\text{Hybrid v50\_2}} &= 1.362e+26\%, \quad \sigma_{\text{Hybrid v50\_2}} = 4.717e+26\% \\
s_{\text{pooled}} &= \sqrt{\frac{\sigma_{\text{Neural Network}}^2 + \sigma_{\text{Hybrid v50\_2}}^2}{2}} = 3.336e+26\% \\
d &= \frac{8899923.4 - 1.362e+26}{3.336e+26} = -0.41
\end{align}

\paragraph{Conservative Estimate:}
Using only Neural Network standard deviation (most conservative approach):
\begin{equation}
d_{\text{conservative}} = \frac{8899923.4 - 1.362e+26}{33577002.8} = -4.056e+18 \quad \text{(interpret magnitude per Cohen 1988 conventions)}
\end{equation}

\paragraph{Alternative Calculation (Glass's $\Delta$):}
Using only the control group (Neural Network) standard deviation:
\begin{equation}
\Delta = \frac{8899923.4 - 1.362e+26}{33577002.8} = -4.056e+18
\end{equation}

\paragraph{Interpretation caveat.} The Hybrid v50\_2 mean and SD are dominated by a single equation whose
extrapolation error is $1.634\times10^{27}\,\%$, so the mean-based effect sizes above (including the
$-4.056\times10^{18}$ ``conservative'' value, which is $(\mu_{NN}-\mu_{H})/\sigma_{NN}$) are not interpretable.
The medians (2415.2\,\% against 13283.9\,\%) and the rank-based $U$ test are the meaningful quantities;
the test does not reject $H_0$. (Recomputed 2026-09-29 from the canonical \texttt{hypatiax/data} files; the previous
\texttt{data\_old}-based block gave $n=11$, $U=68.0$, $p=0.234$.)

% Source: exp1_five_results.json (hypatiax/data/results/five_systems/exp1_five/exp1_five_results.json)

\subsection*{Experiment 2 (exp2\_five, 10-domain suite)}

\textit{No statistical test is reported for exp2\_five: the generator requires both Hybrid v50\_2 and Neural Network to have at least one successful, finite extrapolation-error measurement each, and exp2\_five currently has zero for both (see the corresponding tables in the ``Five-System Comparative Analysis'' section of the main paper).}

\subsection{Interpolation Performance Comparison}

\paragraph{Kruskal-Wallis H Test:}
\textbf{[Removed: a Kruskal-Wallis statistic ($H = 42.3$, $p < 0.001$) and the
``significant differences across methods'' interpretation previously stated here.
No data or script behind them was located, and the per-system sample sizes in the
table below are unequal and unverified. Withdrawn pending recomputation from
source data.]}

\paragraph{System Comparisons:}
\begin{table}[H]
\centering
\caption{Interpolation Performance Statistics by System.
  \textbf{[Unverified: no source data located for these summary statistics; not
  compared against result files (see the Post-Submission Report). Cross-check against the
  main paper's regenerated Table~\texttt{tab:five\_systems\_performance} (Core-15): System~3 has
  $n=15$ there, not $30$; Pure PySR is absent there; and the Hybrid v50\_2 and Neural Network
  means ($0.931$, $0.940$) differ from the regenerated train-$R^2$ means ($0.997$, $0.982$).
  These rows appear to predate the regeneration.]}}
\label{tab:interpolation_stats}
\begin{tabular}{lccccc}
\toprule
\textbf{System} & \textbf{n} & \textbf{Mean $R^2$} & \textbf{Median $R^2$} & \textbf{SD} & \textbf{Range} \\
\midrule
System 3 (LLM+Fallback) & 30 & 1.000 & 1.000 & 0.000 & [1.00, 1.00] \\
Hybrid v50\_2 & 15 & 0.931 & 1.000 & 0.147 & [0.58, 1.00] \\
Neural Network & 15 & 0.940 & 0.952 & 0.082 & [0.78, 1.00] \\
Pure PySR & 18 & 0.982 & 0.998 & 0.045 & [0.88, 1.00] \\
\bottomrule
\end{tabular}
\end{table}

\paragraph{Key Observation:} System 3 achieves R$^{2}$ = 1.000 on interpolation (n=30) but was not designed for extrapolation testing. Hybrid v50\_2's median (1.000) exceeds its mean (0.931) due to outlier presence, making median more representative of typical performance.

\subsection{Domain-Specific Success Rates}

Success rate defined as R$^{2}$ $> 0.95$:

\begin{table}[H]
\centering
\caption{Success Rate by Domain and Method.
  \textbf{[Unverified: not regenerated from result files; the percentages (multiples
  of 33\%) imply about three tasks per domain. Only the column averages were
  checked.]}}
\label{tab:domain_success_detailed}
\begin{tabular}{lccccc}
\toprule
\textbf{Method} & \textbf{Physics} & \textbf{Chemistry} & \textbf{Biology} & \textbf{DeFi} & \textbf{Economics} \\
\midrule
Hybrid v50\_2 & 100\% & 100\% & 67\% & 100\% & 67\% \\
Pure PySR & 100\% & 100\% & 100\% & 67\% & 67\% \\
LLM-Guided PySR & 100\% & 67\% & 67\% & 67\% & 67\% \\
Neural Network & 100\% & 67\% & 33\% & 67\% & 33\% \\
Pure LLM & 100\% & 33\% & 0\% & 33\% & 33\% \\
\midrule
\textbf{Avg across methods} & 100\% & 73\% & 53\% & 67\% & 53\% \\
\bottomrule
\end{tabular}
\end{table}

\paragraph{Correlation Analysis:} We did not measure equation prevalence, so the reproductive-behaviour hypothesis of the main paper is not tested by a correlation here.

\clearpage
\subsection{Validation Layer Effectiveness}

Error detection breakdown across four-layer validation framework:

\begin{table}[H]
\centering
\caption{Validation Layer Error Detection Statistics.
  \textbf{[Unverified: error counts not regenerated from result files; only the
  percentages and cumulative column were checked for internal consistency.]}}
\label{tab:validation_stats}
\begin{tabular}{@{}lcccp{3cm}@{}}
\toprule
\textbf{Layer} & \textbf{Errors} & \textbf{\%} & \textbf{Cum.} & \textbf{Example} \\
\midrule
1. Dimensional Analysis & 18 & 12 & 12 & Energy = mass \\
2. Physics Constraints & 27 & 18 & 30 & Negative rate \\
3. Numerical Stability & 35 & 23 & 53 & Overflow \\
4. Prediction Accuracy & 70 & 47 & 100 & Wrong result \\
\midrule
\textbf{Total} & \textbf{150} & \textbf{100} & --- & --- \\
\bottomrule
\end{tabular}
\end{table}

\paragraph{Key Finding:} Single-layer validation (prediction accuracy only) would miss 53\% of errors caught by earlier layers on this benchmark, indicating that cascading validation contributed substantially to error detection for LLM-generated outputs.

\clearpage
% ============================================================================
% APPENDIX E: DISCOVERED EXPRESSION EXAMPLES
% ============================================================================


\section{Discovered Expression Examples}
\label{app:expressions}

\subsection{Perfect Recoveries}

\subsubsection{Kinetic Energy}

\begin{example}[Kinetic Energy Discovery]
\textbf{Ground truth}: $E = \frac{1}{2}mv^2$

\textbf{HypatiaX v50\_2 discovered}:
\begin{equation}
E = ((v \times m) \times v) \times 0.5 = 0.5mv^2
\end{equation}

\textbf{Analysis}: Exact recovery of ground truth. Discovered expression is algebraically identical.

\textbf{Interpolation}: $R^2 = 1.0000$ on training data (v $\in$ [1, 10] m/s)

\textbf{Extrapolation}: 0.0\% error at v = 50 m/s (5$\times$ training max)

\textbf{Physical interpretation}: Correctly captures quadratic dependence on velocity and linear dependence on mass.
\end{example}

\subsubsection{Ideal Gas Law}

\begin{example}[Ideal Gas Law Discovery]
\textbf{Ground truth}: $P = nRT/V$, where $R = 8.314$\,J/(mol$\cdot$K)

\textbf{HypatiaX v50\_2 discovered}:
\begin{equation}
P = n \times (T \times (8.314 / V))
\end{equation}

\textbf{Analysis}: Exact recovery including the universal gas constant $R = 8.314$ to full precision.

\textbf{Interpolation}: $R^2 = 1.0000$ on training data

\textbf{Extrapolation}: 0.0\% error across all pressure/temperature/volume ranges tested (up to 5$\times$ training range)

\textbf{Physical interpretation}: On this test case, symbolic regression recovered a constant consistent with the gravitational constant from data alone, without prior knowledge of its value\citep{lacava2021contemporary}.
\end{example}

\subsection{Near-Perfect Recoveries}

The Michaelis-Menten and Impermanent Loss equations were recovered with constant error $< 10^{-5}$, achieving $R^2 = 1.0000$ and near-zero extrapolation error (see the discovered-expression examples above).\textbf{[EDITORIAL NOTE: original text pointed to ``Listing~\texttt{\textbackslash ref\{sec:results\}}''; \texttt{sec:results} is the main paper's entire Results section, not a listing, so this appears to be a leftover/incorrect reference. Removed pending confirmation of intended target.]}

\subsection{Failure Case Analysis}

\subsubsection{Gravitational Force (Both Methods Failed)}

\begin{example}[Gravitational Force Failure]
\textbf{Ground truth}: $F = Gm_1m_2/r^2$, where $G = 6.674 \times 10^{-11}$\,N$\cdot$m$^2$/kg$^2$

\textbf{Neural Network Result}:
\begin{itemize}
\item \textbf{Interpolation}: $R^2 = 0.21$ (poor fit; relative error $> 100\%$, meeting the threshold in the main paper's \emph{Interpolation--Extrapolation Decoupling} definition)
\item \textbf{Extrapolation}: Not attempted due to interpolation failure
\item \textbf{Issue}: Gradient vanishing due to $10^{-11}$ scale
\item \textbf{Attempted solution}: Log-transform $\to$ marginal improvement ($R^2 = 0.35$)
\item \textbf{Root cause}: Optimization landscapes too flat at this scale for gradient descent
\end{itemize}

\textbf{HypatiaX v50\_2 Result}:
\begin{itemize}
\item \textbf{Interpolation}: $R^2 = -0.03$ (worse than baseline)
\item \textbf{Extrapolation}: Not attempted due to interpolation failure
\item \textbf{Discovered form}: $F = 1.23 \times 10^{13} \frac{m_1 m_2}{r^2}$ (correct functional form, wrong constant by $2 \times 10^{23}\times$)
\item \textbf{Issue}: PySR constant search range [$10^{-3}$, $10^3$] missed $G \sim 10^{-11}$
\item \textbf{Attempted solution}: Extend range to [$10^{-15}$, $10^{15}$] $\to$ no convergence in 50 iterations
\item \textbf{Root cause}: Search space too large for evolutionary algorithm to explore efficiently
\end{itemize}

\textbf{Recommended Solution}:
\begin{enumerate}
\item \textbf{Dimensional analysis preprocessing}: Identify required constant scales before symbolic regression
\item \textbf{Logarithmic transformation}: Convert $F = Gm_1m_2/r^2$ to $\log F = \log G + \log m_1 + \log m_2 - 2\log r$, making the problem linear in log-space
\item \textbf{Adaptive constant ranges}: Dynamically adjust search based on data magnitude
\item \textbf{Multi-scale normalization}: Normalize inputs/outputs to [0.1, 10] range before discovery, then rescale
\end{enumerate}

\textbf{Lesson learned}: Very small constants ($< 10^{-6}$) require special handling in both neural and symbolic methods. This represents 1/15 test cases (6.7\% failure rate) and is a known limitation addressed in Section~\ref{suppb:limitations}.
\end{example}

\subsubsection{Henderson-Hasselbalch (Pure LLM Failure)}

\begin{example}[Henderson-Hasselbalch pH Calculation]
\textbf{Ground truth}: $\text{pH} = \text{pKa} + \log_{10}\left(\frac{[\text{A}^-]}{[\text{HA}]}\right)$

\textbf{Pure LLM Generated}:
\begin{equation}
\text{pH} = \text{pKa} \times \frac{[\text{A}^-]}{[\text{HA}]}
\end{equation}

\textbf{Failure Analysis}:
\begin{itemize}
\item \textbf{R$^{2}$ score}: $-12.3$ (predictions anticorrelated with ground truth)
\item \textbf{Errors}:
  \begin{enumerate}
  \item Multiplication instead of addition
  \item Missing logarithm entirely
  \item Plausible but fundamentally wrong
  \end{enumerate}
\item \textbf{Root cause}: LLM training data likely contains more qualitative discussions of pH (''pH increases when conjugate base increases") than quantitative formulas
\item \textbf{Pattern matching failure}: LLM defaulted to simpler multiplicative relationship rather than logarithmic form
\end{itemize}

\textbf{Symbolic Rescue}:
PySR correctly discovered logarithmic relationship:
\begin{equation}
\text{pH} = \text{pKa} + \log_{10}\left(\frac{[\text{A}^-]}{[\text{HA}]}\right)
\end{equation}
Achieving $R^2 = 0.998$, demonstrating the value of hybrid architecture.

\textbf{Lesson learned}: LLM performance declined on this specialised equation, which is consistent with the main paper's \emph{LLM Reproductive Behaviour} finding---generation probability is theoretically predicted to be lower for expressions that are structurally distant from the training corpus\citep{kambhampati2024llms}.
\end{example}

\subsection{Complete Discovered Expressions Code Listing}

\begin{lstlisting}[caption=Discovered Symbolic Expressions (7 Perfect Recoveries), label=lst:discovered, language=Python]
# Kinetic Energy (PERFECT)
discovered: ((v * m) * v) * 0.5
ground_truth: 0.5 * m * v**2
R2: 1.0000, extrapolation_error: 0.0%

# Rate Law (PERFECT)
discovered: ((A_conc * 0.5) * A_conc) * B_conc
ground_truth: 0.5 * A_conc**2 * B_conc
R2: 1.0000, extrapolation_error: 0.0%

# Ideal Gas Law (PERFECT)
discovered: n * (T * (8.314 / V))
ground_truth: n * 8.314 * T / V
R2: 1.0000, extrapolation_error: 0.0%

# Logistic Growth (PERFECT - algebraically equivalent)
discovered: ((N * -0.0003) + 0.3) * N
ground_truth: 0.3 * N * (1 - N/1000)
expanded: 0.3*N - 0.0003*N^2 (identical)
R2: 1.0000, extrapolation_error: 0.0%

# Price Impact (PERFECT)
discovered: swap / (swap + reserve)
ground_truth: dx / (x + dx)
R2: 1.0000, extrapolation_error: 0.0%

# Michaelis-Menten (NEAR-PERFECT)
discovered: (S / (S + 10.000007)) * 50.000008
ground_truth: 50 * S / (10 + S)
constant_error: <1e-5
R2: 1.0000, extrapolation_error: 0.0%

# Impermanent Loss (NEAR-PERFECT)
discovered: (2 * sqrt(r)) / (1 + r) - 1.000023
ground_truth: 2 * sqrt(r) / (1 + r) - 1
constant_error: 2.3e-5
R2: 1.0000, extrapolation_error: 0.0%
\end{lstlisting}

% ============================================================================
% APPENDIX F: REPRODUCIBILITY STATEMENT (CORRECTED 2026-03-21)
% ============================================================================


\section{Supplementary Materials}
\label{app:supplementary}

The following supplementary materials are available:

\begin{itemize}
\item \textbf{S1: Pure LLM Baseline Results} (\texttt{S1\_pure\_llm\_results.csv}) --- Complete results from Pure LLM baseline experiments.

\item \textbf{S2: Extrapolation Test Data} (\texttt{S2\_extra\-polation\_data.csv}) --- Training R$^2$ scores and extrapolation errors at 1.2$\times$, 2$\times$, and 5$\times$ training ranges.

\item \textbf{S3: Validation Logs} (\texttt{S3\_validation\_logs.json}) --- Complete validation reports for all 131 tests showing the four-layer validation framework in action. Each entry includes dimensional analysis, physics constraints, numerical stability checks, prediction accuracy metrics, timing data, and failure mode classification. Includes validation layer rejection rates (12\%, 18\%, 23\%, 47\%).

\item \textbf{S4: Discovered Symbolic Expressions} (\texttt{S4\_discovered\_expressions.\{txt,json\}}) --- All 131 discovered symbolic expressions in SymPy-compatible format, organized by domain (chemistry, biology, physics, DeFi, finance) with complete metadata including ground truth equations, performance metrics (R$^2$, timing, iterations), and algebraic equivalence verification.

\item \textbf{S5: Timing Breakdown} (\texttt{S5\_timing\_breakdown.csv}) --- Detailed computational cost analysis showing component-level timing: LLM initialization (0-39\%), symbolic search (50-91\%), and validation (7-11\%). Demonstrates the 4.3× speedup from LLM-guided initialization versus pure symbolic methods.

\item \textbf{S6: Domain-Specific Analysis} (\texttt{S6\_domain\_analysis.xlsx}) --- Statistical analysis by scientific domain showing success rates: Biology 100\%, Chemistry 100\%, Physics 100\%, DeFi 100\%, Finance 77.8\%. Includes per-equation breakdowns and domain-specific constraint validation results.

\item \textbf{S7: Failure Mode Taxonomy} (\texttt{S7\_failure\_mode\_taxonomy.pdf}) --- Comprehensive taxonomy of 23 distinct failure modes categorized by method (Pure LLM: 8 modes, Neural Network: 3 modes, Symbolic: 5 modes, Hybrid: 2 modes, Integration: 5 modes), severity (8 Critical, 4 High, 9 Medium, 2 Low), and detection layer. Includes detailed examples of silent semantic errors, large-scale extrapolation errors, distributional reasoning failures, and API misuse.

\item \textbf{S8: Figure Source Data} (\texttt{S8\_figures\_source\_data/}) --- Source data files for all figures in the paper:
\begin{itemize}
\item \texttt{figure1\_arrhenius\_extrapolation.csv} --- Arrhenius equation extrapolation comparison
\item \texttt{figure2\_domain\_comparison.csv} --- Success rates by scientific domain
\item \texttt{figure3\_validation\_breakdown.csv} --- Four-layer validation effectiveness
\item \texttt{figure4b\_domain\_breakdown.csv} --- Per-equation performance details
\item \texttt{figure5\_method\_comparison.csv} --- Overall method comparison
\item \texttt{figure\_5systems\_comparison.csv} --- Extrapolation data for the five-system comparison; see the ``Five-System Comparative Analysis'' section of the main paper for the regenerated tables (see Issue 4 correction note above)
\item \texttt{figure\_benchmark\_comparison.csv} --- Core 15 benchmark results
\item \texttt{figure\_timing\_comparison.csv} --- Speed-accuracy tradeoff analysis
\end{itemize}
\end{itemize}

All supplementary materials are provided in both synthetic versions (matching reported statistics exactly) and real experimental versions (direct from JSON logs) for maximum reproducibility and transparency. Complete documentation and usage examples are provided in accompanying README files.

\subsection{Supplementary Data Files}

Post-processed results, tables and the paper sources are available in the
post-processing and pre-submission repository:
\begin{center}
{\small\url{https://github.com/sednabcn/LLM-HypatiaX-DEV}}
\end{center}
The raw benchmark runs are produced and archived in the reproduction
repository, \url{https://github.com/sednabcn/LLM-HypatiaX-REPRO}.

\begin{table}[H]
\centering
\caption{Supplementary Data Files}
\label{tab:supplementary_files}
\small
\begin{tabular}{@{}lp{8cm}@{}}
\toprule
\textbf{File} & \textbf{Contents} \\
\midrule
\texttt{S1\_pure\_llm\_results.csv} & Complete Pure LLM test results (n=40) with formulas, R$^{2}$ scores, failure modes \\
\texttt{S2\_extrapolation\_data.csv} & Extrapolation test data ($n=27$) at $1.2\times$, $2\times$, $5\times$ training ranges \\
\texttt{S3\_validation\_logs.json} & Full validation reports for all 131 tests \\
\texttt{S4\_discovered\_expressions.txt} & All discovered symbolic expressions in SymPy format \\
\texttt{S5\_timing\_breakdown.csv} & Detailed timing data by component and method \\
\texttt{S6\_domain\_analysis.xlsx} & Per-domain performance with statistical tests \\
\texttt{S7\_failure\_taxonomy.pdf} & Complete taxonomy of 23 failure modes with examples \\
\texttt{S8\_figures\_source\_data/} & Raw data for all 8 figures in CSV format \\
\bottomrule
\end{tabular}
\end{table}

\subsection{Reproducibility Materials}

Complete step-by-step tutorials for reproducing all experiments are available at:
\url{https://sednabcn.github.io/ai-llm-blog/tutorials/hypatiax/}

\begin{itemize}
\item \href{https://sednabcn.github.io/ai-llm-blog/tutorials/hypatiax/setup/}{Tutorial 1: Environment Setup and First Discovery} (15 min)
\item \href{https://sednabcn.github.io/ai-llm-blog/tutorials/hypatiax/experiments/}{Tutorial 2: Running Benchmark Experiments} (45 min active + 3--8 hours compute)
\item \href{https://sednabcn.github.io/ai-llm-blog/tutorials/hypatiax/analysis/}{Tutorial 3: Statistical Analysis and Publication Figures} (30 min)
\item \href{https://sednabcn.github.io/ai-llm-blog/tutorials/hypatiax/extensions/}{Tutorial 4: Custom Applications and Extensions} (45 min)
\end{itemize}

All source code, experimental protocols and data generation scripts are available at:
\url{https://github.com/sednabcn/LLM-HypatiaX-REPRO}
(reproduction repository: all tests and benchmarks are run from there;
\url{https://github.com/sednabcn/LLM-HypatiaX-DEV} holds post-processing and pre-submission material only).

\subsection{Reproducibility Artifacts}

\paragraph{Docker Images:}
\begin{itemize}
\item \texttt{hypatiax:v1.0} - Full environment with all dependencies
\item \texttt{hypatiax:minimal} - Core symbolic regression only (no LLM)
\item \texttt{hypatiax:gpu} - GPU-accelerated version for neural network comparison
\end{itemize}

\paragraph{Source Code:}
All Python scripts and experimental protocols are available at:
\begin{center}
{\small\url{https://github.com/sednabcn/LLM-HypatiaX-REPRO/tree/master/hypatiax}}
\end{center}


\paragraph{Jupyter Notebooks:}

The notebooks that do exist are the six paper-integrity audit notebooks
\texttt{NB-01} to \texttt{NB-06} (citation, cross-reference, section numbering,
numerical consistency, figure dependency, and code-quality checks), available at:
\begin{center}
{\footnotesize\url{https://github.com/sednabcn/LLM-HypatiaX-DEV/tree/master/notebooks}}
\end{center}


\paragraph{Step-by-Step Tutorials:}
\begin{itemize}
\item \href{https://sednabcn.github.io/ai-llm-blog/tutorials/hypatiax/setup/}{Tutorial 1: Environment Setup and First Discovery} (15 min)
\item \href{https://sednabcn.github.io/ai-llm-blog/tutorials/hypatiax/experiments/}{Tutorial 2: Running Benchmark Experiments} (45 min active + 3--8 hours compute)
\item \href{https://sednabcn.github.io/ai-llm-blog/tutorials/hypatiax/analysis/}{Tutorial 3: Statistical Analysis and Publication Figures} (30 min)
\item \href{https://sednabcn.github.io/ai-llm-blog/tutorials/hypatiax/extensions/}{Tutorial 4: Custom Applications and Extensions} (45 min)
\end{itemize}
Available at: \url{https://sednabcn.github.io/ai-llm-blog/tutorials/hypatiax/}


\subsection{Additional Statistical Tests}

\begin{table}[H]
\centering
\caption{Additional Statistical Tests.
  \textbf{[Unverified, manually entered: no script or data located that reproduces
  these test statistics or $p$-values; do not cite as computed results.]}}
\label{tab:additional_stats}
{\small\begin{tabular}{@{}lccl@{}}
\toprule
\textbf{Comparison} & \textbf{Test} & \textbf{Result} & \textbf{Interpretation} \\
\midrule
Hybrid v50\_2 vs Pure PySR & Mann-Whitney\textsuperscript{$\dagger$} & $U = 89$, $p = 0.18$ & No significant difference \\
System 3 vs Hybrid v50\_2 & Mann-Whitney\textsuperscript{$\dagger$} & $U = 156$, $p = 0.03$ & System 3 better (interp.) \\
Classical vs DeFi (LLM) & Fisher's Exact\textsuperscript{$\dagger$} & $p < 0.001$ & Domain effect significant \\
Run-to-run variability & Friedman\textsuperscript{$\dagger$} & $\chi^2 = 34.2$, $p < 0.001$ & High variability \\
\bottomrule
\end{tabular}}

\smallskip
{\footnotesize\textbf{[$^\dagger$Not reproducible.]} No raw per-equation values or script
reproducing these four rows were located. (i)~The two Mann-Whitney rows disagree with this
file's own sample sizes (\Cref{tab:interpolation_stats}: $n=15$, $18$, $30$): with those
sizes, $U=89$ gives $p\approx0.10$ (two-sided, normal approximation, no tie correction),
not $0.18$, and $U=156$ gives $p\approx0.10$, not $0.03$. (ii)~Pure PySR is not one of the
systems in the main paper's regenerated Core-15 five-system tables. (iii)~The Fisher and
Friedman rows give no counts, degrees of freedom or data, so they cannot be checked at all.
The only test in this document recomputed from source files is the Hybrid v50\_2 vs.\ Neural
Network extrapolation test in Appendix~\ref{app:statistical_tests}.}
\end{table}

\paragraph{Power Analysis:}
\begin{itemize}
\item Domain comparisons (n=3 per domain): Power = 0.65 for d > 0.8 \textbf{[Unverified: no calculation or source located.]}
\item Method comparisons (n=15 per method): Power = 0.95 for d > 0.6 \textbf{[Unverified: no calculation or source located; ``n=15 per method'' also disagrees with the sample sizes $15$/$15$/$18$/$30$ in \Cref{tab:interpolation_stats}.]}
\end{itemize}

Recommendation: Future work should increase per-domain sample sizes to n $\geq$ 10 for robust inference on secondary effects.

% ============================================================================
% APPENDIX H: NOISE SENSITIVITY ANALYSIS
% ============================================================================


\section{Noise Sensitivity Analysis}
\label{app:noise_sensitivity}

This section reports solve counts for the five noise levels present in the
repository's \texttt{protocol\_core} result files ($\sigma\leq1\%$). No results are available for higher noise levels, so no claim is made about a noise threshold.

\subsection{Protocol}

Five \texttt{protocol\_core} runs (one noiseless, four noisy; $200$ samples each)
cover $\sigma\in\{0,\,0.05,\,0.1,\,0.5,\,1\}\%$, read from each file's
\texttt{protocol.noise\_level}, with 30 equations per level (150 records per
system). A record counts as solved if $\Rsq\geq$ the threshold recorded in its
own file's protocol block (\texttt{protocol.threshold}); the
\texttt{success} flag is not used (M3 sets it to true whenever no exception is
raised).

\subsection{Results}

\begin{table}[H]
\centering
\caption{Solved counts per noise level ($\Rsq\geq$ the run's own recorded
  threshold; $n=30$ per cell). Regenerated from the five \texttt{protocol\_core}
  files with \texttt{generate\_tables.py} (\texttt{noise\_sensitivity.tex});
  columns give $\sigma$ in percent. \textbf{Thresholds differ per run}
  ($0.999999$, $0.95$, $0.95$, $0.995$, $0.99$ for $\sigma=0$, $0.05$,
  $0.1$, $0.5$, $1\%$; these are the thresholds written into each
  \texttt{protocol\_core} file, and differ from the $0.995$ that
  \Cref{tab:rr_noise} applies at $\sigma=0.05\%$ and $0.1\%$), so counts are \emph{not} comparable across noise
  levels, and the pooled column adds counts taken under different thresholds.
  \textbf{Source of the $\sigma=0$ column:} all six cells come from the noiseless run of the noise sweep, \texttt{protocol\_core\_noiseless\_20260812\_170344.json}; no cell is taken from \texttt{benchmark\_results.json} (see the tablenote on \Cref{tab:overall}).}
\label{tab:noise_sensitivity}
\small
\setlength{\tabcolsep}{4pt}
\begin{tabular}{@{}l c c c c c r@{}}
\toprule
\textbf{Method, $\sigma$ (\%)} & \textbf{0} & \textbf{0.05} & \textbf{0.1} & \textbf{0.5} & \textbf{1} & \textbf{Pooled}\\
\midrule
\EHD{} (M3)            & 23/30 & 30/30 & 30/30 & 30/30 & 30/30 & 143/150 (95.3\%)\\
HybridDiscovery v50\_2 & 21/30 & 25/30 & 26/30 & 23/30 & 26/30 & 121/150 (80.7\%)\\
\HSL{} (M4)            & 30/30 & 30/30 & 30/30 & 30/30 & 30/30 & 150/150 (100.0\%)\\
ImprovedNN             &  0/30 & 29/30 & 29/30 & 27/30 & 29/30 & 114/150 (76.0\%)\\
PureLLM                & 30/30 & 29/30 & 29/30 & 30/30 & 30/30 & 148/150 (98.7\%)\\
SymbolicEngineWithLLM  & 28/30 & 30/30 & 30/30 & 29/30 & 29/30 & 146/150 (97.3\%)\\
\bottomrule
\end{tabular}
\end{table}

\begin{table}[H]
\centering
\caption{Solved counts at a single fixed threshold, $\Rsq\geq0.99$, for every
  run ($n=30$ per cell). Because the threshold is the same at all noise levels,
  counts are comparable across levels and the pooled column is valid.
  Regenerated with \texttt{generate\_tables.py --fixed-threshold 0.99}
  (\texttt{noise\_sensitivity\_fixed\_0p99.tex}); columns give $\sigma$ in percent. The $\sigma=0$ column is read from \texttt{protocol\_core\_noiseless\_20260812\_170344.json}, as in \Cref{tab:noise_sensitivity}.}
\label{tab:noise_sensitivity_fixed}
\small
\setlength{\tabcolsep}{4pt}
\begin{tabular}{@{}l c c c c c r@{}}
\toprule
\textbf{Method, $\sigma$ (\%)} & \textbf{0} & \textbf{0.05} & \textbf{0.1} & \textbf{0.5} & \textbf{1} & \textbf{Pooled}\\
\midrule
\EHD{} (M3)            & 30/30 & 30/30 & 30/30 & 30/30 & 30/30 & 150/150 (100.0\%)\\
HybridDiscovery v50\_2 & 30/30 & 24/30 & 26/30 & 24/30 & 26/30 & 130/150 (86.7\%)\\
\HSL{} (M4)            & 30/30 & 30/30 & 30/30 & 30/30 & 30/30 & 150/150 (100.0\%)\\
ImprovedNN             & 29/30 & 27/30 & 28/30 & 29/30 & 29/30 & 142/150 (94.7\%)\\
PureLLM                & 30/30 & 29/30 & 29/30 & 30/30 & 30/30 & 148/150 (98.7\%)\\
SymbolicEngineWithLLM  & 30/30 & 30/30 & 30/30 & 30/30 & 29/30 & 149/150 (99.3\%)\\
\bottomrule
\end{tabular}
\end{table}

\subsection{Key Findings}

\paragraph{Ranking.} At the fixed $\Rsq\geq0.99$ threshold, M3 and M4 both solve
150/150, ahead of SymbolicEngineWithLLM (149), PureLLM (148), ImprovedNN (142)
and HybridDiscovery v50\_2 (130). M3's apparent shortfall in the per-run table
(23/30 at $\sigma=0$) comes from the strict $0.999999$ threshold used for that
run, not from a loss of accuracy. At that strict threshold M3 solves 23/30
noiseless equations while M4 solves 30/30, so M4 is the only system that is exact
under the strict criterion. At $0.99$ the two systems cannot be separated.

\paragraph{Threshold confound.} The per-run table mixes thresholds from $0.95$ to
$0.999999$. For example ImprovedNN scores 0/30 at $\sigma=0$ under $0.999999$ but
29/30 under $0.99$; the jump in the per-run table is the threshold, not the
data. Cross-level comparisons should use the fixed-threshold table. The value
$0.99$ is permissive, so it does not measure exact recovery.

\paragraph{Scope.} Only $\sigma\leq1\%$ is covered, and at
$\Rsq\geq0.99$ essentially all systems saturate, so these data do not show where
noise starts to hurt. No failure-mode, domain-variability, or noise-threshold claim is supported by
these data.

% ============================================================================
% APPENDIX I: EXTENDED DISCUSSION
% ============================================================================




\end{document}
